%!TEX program = pdflatex
% Mathematics article -- amsart. Target: J. Phys. A: Math. Theor. (regular Paper).
\documentclass[11pt,reqno]{amsart}

\usepackage[margin=1in]{geometry}

\usepackage{amssymb,mathtools}
\usepackage{graphicx}
\usepackage{bm}
\usepackage{xcolor}
\usepackage{array,booktabs,tabularx,ragged2e}
\usepackage{cite}
\usepackage{hyperref}
\hypersetup{colorlinks=true,
            linkcolor=blue!55!black,
            citecolor=blue!55!black,
            urlcolor=blue!55!black}
\usepackage[capitalize,nameinlink]{cleveref}

\theoremstyle{plain}
\newtheorem{theorem}{Theorem}[section]
\newtheorem{lemma}[theorem]{Lemma}
\newtheorem{proposition}[theorem]{Proposition}
\newtheorem{corollary}[theorem]{Corollary}
\newtheorem{conjecture}[theorem]{Conjecture}

\theoremstyle{definition}
\newtheorem{definition}[theorem]{Definition}
\newtheorem{example}[theorem]{Example}

\theoremstyle{remark}
\newtheorem{remark}[theorem]{Remark}

\numberwithin{equation}{section}

\setlength{\emergencystretch}{3em}

\newcommand{\Mat}{\operatorname{Mat}}
\newcommand{\Tr}{\operatorname{Tr}}
\providecommand{\rank}{\operatorname{rank}}
\providecommand{\tr}{\operatorname{tr}}
\providecommand{\diag}{\operatorname{diag}}
\newcommand{\C}{\mathbb C}
\newcommand{\Id}{I_2}
\renewcommand{\arraystretch}{1.16}

\title[Reflection matrices for six difference-form $R$-matrices]
      {Complete classification of reflection matrices for six
       difference-form spin-chain $R$-matrices}

\author{Zhiyuan Yao}
\address{Lanzhou Center for Theoretical Physics, Key Laboratory of Theoretical
  Physics of Gansu Province, Key Laboratory of Quantum Theory and Applications
  of MoE, Gansu Provincial Research Center for Basic Disciplines of Quantum
  Physics, Lanzhou University, Lanzhou 730000, China}
\email{yaozy@lzu.edu.cn}

\date{\today}

\subjclass[2020]{Primary 16T25; Secondary 81R12, 82B23, 13P10}
\keywords{reflection equation, boundary Yang--Baxter equation, integrable spin
  chains, open chains, $K$-matrix, constructible classification}

\begin{document}

\begin{abstract}
An integrable quantum spin chain acquires an integrable open end when a
reflection matrix solves the boundary Yang--Baxter equation of the chain's bulk
$R$-matrix. For the six regular difference-form $4\times4$ $R$-matrices that
fall outside the eight-vertex list no such matrix has been known: the bulk
classification that produced them concerns closed chains, and most of the six
are not $PT$ symmetric, so even the ordering of the reflection equation is part
of the problem. We classify, for each of the six and at every point of its full
complex parameter variety, every single-valued meromorphic $2\times2$ reflection
matrix that is holomorphic at the origin and proportional to the identity there.
The classification starts from four unrestricted meromorphic entries: one
differentiation of the reflection equation at the origin produces a homogeneous
linear system of sixteen equations in those entries, and an exact stratification
of each parameter variety by the rank of that system splits it into finitely
many cells carrying either a unique kernel or a family of arbitrary functions
fixed by a centralizer. The tables include every removable limit, singular
stratum, arbitrary-function fiber and overlap, and every family continues to a
single-valued meromorphic function on the plane. Four of the six classes admit
no crossing datum whatsoever, because the first-leg partial transpose of their
bulk matrices has a zero row and a zero column; the remaining two admit a unique
one exactly where the rate parameter is nonzero, but the left boundary that
datum prescribes yields commuting double-row transfer matrices for every
independently chosen pair of boundaries only on the rational cells. Fixing the
left boundary to the identity instead removes that restriction for Class~6
entirely: every classified right boundary, at every point of the Class-6
parameter plane and for every chain length, gives a commuting double-row family
whose logarithmic derivative at the origin is twice the bulk Hamiltonian plus a
one-site endpoint term set by the derivative of the boundary matrix.
\end{abstract}

\maketitle

\section{Introduction}
\label{sec:intro}

A quantum spin chain is integrable when it carries a tower of mutually commuting
conserved charges, one of which is its Hamiltonian, and all of which are
generated by a single object: an invertible matrix depending on a complex
spectral parameter and solving the Yang--Baxter equation, the identity
guaranteeing that multi-particle scattering factorizes into two-particle events
\cite{Yang_1967se,Baxter_1972pf,Takhtadzhan_1979tq}. Because everything else
follows from that one matrix, the natural question at a fixed local dimension is
not which models are integrable but how many there are, and the answer has been
pursued by direct classification since the constant case was settled
\cite{Kulish_1982so,Hietarinta_1993st,Vieira_2018sa,Retore_2022it}. For a chain
of spins one-half, with the bulk matrix depending on its two rapidities only
through their difference and reducing at coincident rapidities to the exchange
of the two sites, the classification is now complete: fourteen independent
solutions, of which eight are the familiar vertex models and six lie outside
that list, several of them with nilpotent or non-diagonalizable Hamiltonian
densities \cite{deLeeuw2019}. Those six have since been absorbed into a
completed classification of all regular $4\times4$ solutions, which drops the
difference-form assumption and adds four further models
\cite{Leeuw_2021yb,corcoran2024}. The fifth of them has been analysed as a closed
chain whose transfer matrix and Hamiltonian are both non-diagonalizable, so that
the algebraic Bethe ansatz does not reach the whole state space
\cite{nieto2024}, and identified as a Jordanian Drinfeld twist of the Heisenberg
chain whose continuum limit is a non-unitary deformation of the
Landau--Lifshitz model, with the sixth treated the same way \cite{deLeeuw2026}.
The classification method itself has been carried to electron chains with a
four-dimensional local space \cite{Leeuw_2020ni}, and to the nearest-neighbour
interactions underlying deformations of the $AdS_5\times S^5$ worldsheet
model \cite{Leeuw_2020cn}. The commuting charges these matrices generate are
organized by a Yangian, whose representation theory is the structural
background to the whole programme \cite{Loebbert_2016lo}.

Everything in that list concerns a ring. Cutting the ring open at a wall requires
one further piece of algebraic data: a reflection matrix, giving the amplitude
with which an excitation returns from the wall with its rapidity reversed, and
solving the reflection equation --- the statement that two excitations may bounce
off the wall in either order with the same total amplitude
\cite{cherednik1984,Sklyanin_1987bc,sklyanin1988}. That equation is what
preserves the commuting tower after the cut, and its solutions are what turn a
list of bulk models into a catalogue of integrable boundary interactions; in
algebraic language they select the part of the bulk symmetry that survives at the
wall, a coideal subalgebra \cite{kulish1992,Kolb_2009re,Balagovic_2016uk}. For
the standard chains the reflection matrices are known completely
\cite{devega1994}, and for rational reflection matrices that stay invertible as
the spectral parameter grows, the possible residual symmetry algebras have been
derived from the boundary Yang--Baxter equation alone \cite{Gombor_2020ot}. Boundaries have also been found for
individual deformed models, the Jordanian open XXX chain among them
\cite{kulish2010}, and integrable open ends are actively being constructed for
discrete-time circuits built from the same algebraic data
\cite{paletta2025,garciaFernandez2026}. For the six families outside the vertex
list there is nothing: no reflection solution, no crossing relation, and no
open-chain Hamiltonian. The paper that isolated them names precisely this as
future work \cite{deLeeuw2019}.

The nearest thing to a precedent is at a different local dimension. The same
differential approach has been pushed to spin-one vertex models: for the
fifteen-vertex $R$-matrices obeying the ice rule, which are $9\times9$ and not
symmetric, four families of regular bulk matrices were found and their
reflection matrices classified alongside, in the same ordering of the reflection
equation used below \cite{Vieira_2019fv}. That is the shape of result this paper
supplies for the six $4\times4$ difference-form families, which are a different
list of bulk matrices at a different local dimension.

Three features make the gap harder to close than it looks. First, most of the six
bulk matrices are not $PT$ symmetric, and for a bulk matrix that is not, the two
orderings of the bulk factors in the reflection equation are genuinely different
equations with genuinely different solution sets \cite{mezincescu1991}; the
literal equation is therefore part of the problem statement rather than a
convention. Second, even the constant boundary --- the identity matrix, which
looks like no boundary at all --- is a nontrivial condition, because substituting
it leaves an identity between products of bulk matrices at shifted arguments that
holds automatically only when the bulk matrix is symmetric. Third, the standard
route from a right boundary to a left one runs through a crossing relation, the
algebraic form of particle--antiparticle exchange, and nothing guarantees that
these bulk matrices possess one. A list of checked examples would settle none of
this: what is needed is a classification, over the whole of each complex
parameter variety, including the singular strata where a parameter vanishes and a
family degenerates.

This paper supplies that classification and its consequences. We fix the literal
reflection equation, the class of admissible boundaries --- single-valued
meromorphic on the plane, holomorphic at the origin, and there a nonzero multiple
of the identity --- and the equivalence under which boundaries are counted, which
is multiplication by a scalar function together with the constant similarities
and spectral rescalings that preserve the same literal bulk matrix. Every
boundary is then found, at every bulk point, of all six classes. The derivation
uses no ansatz: it begins with four unrestricted meromorphic entries, and a
single differentiation of the reflection equation at the origin converts the
problem into a homogeneous linear system of sixteen equations in those four
entries, whose coefficients are the bulk matrix, its derivative, and the first
derivative of the boundary at the origin. What then decides the answer is the
rank of that system over the field of meromorphic functions, and the rank takes
only three values. Where it is three the kernel is one-dimensional, so a single
normalized matrix is the whole solution set. Where it is two or zero the equation
collapses, after an entire invertible change of variable, to the statement that
the boundary commutes with itself at different rapidities, and the centralizer of
a nonscalar $2\times2$ matrix is two-dimensional, so the solutions are exactly
one arbitrary meromorphic function times one fixed traceless direction. The
partition of each parameter variety into cells of constant rank is computed
exactly, from radicals and associated primes of the ideals generated by the
maximal minors, and it is that computation which carries the exhaustiveness of
the tables.

The consequences for open chains split the six classes sharply. For four of them
no crossing datum exists at any bulk point, and the reason is a single rank
count: the first-leg partial transpose of each of those bulk matrices has a zero
row and a zero column, so it cannot be inverted, while a crossing relation would
force it to be. For the remaining two, crossing exists exactly where the rate
parameter is nonzero and is then unique up to scale; but the left boundary it
prescribes produces commuting double-row transfer matrices for every
independently chosen pair of right boundaries only on the rational cells, and we
give the explicit nonzero commutator entries that rule out the rest.

Requiring every independently chosen pair to work is, however, more than an open
chain needs. One compatible left boundary per right boundary already gives a
commuting family, and for Class~6 there is one, the same for all of them: the
identity. For every point of the Class-6 parameter plane, every classified right
boundary, and every chain length, the double-row transfer matrices built with the
identity on the left commute, and their logarithmic derivative at the origin is
twice the sum of the bulk two-site densities plus a single one-site term at the
end carrying the boundary, with no contribution from the auxiliary end because
both partial traces of the Class-6 density vanish. This is proved on the generic
stratum from the dual reflection identity that the boundary transfer-matrix
theorem requires \cite{Vlaar_2015bt}, extended to the singular stratum by
polynomial continuation in the rate, and settled directly at the point where the
bulk matrix degenerates to the exchange operator. No analogous construction is
claimed for the other five classes.

\Cref{sec:setting} fixes the conventions, prints the six bulk matrices, and
states the question. \Cref{sec:exhaustiveness} derives the differentiated system
and proves the rank trichotomy that makes it exhaustive. \Cref{sec:classification}
gives the classification, class by class. \Cref{sec:quotient} carries out the
fixed-bulk quotient. \Cref{sec:crossing} settles crossing, left boundaries, and
the resulting open chains. \Cref{sec:class6} proves the Class-6 construction.
\Cref{sec:checks} reports the independent checks and says what each does and does
not establish, and \cref{sec:discussion} states the limitations.

\section{Setting}
\label{sec:setting}

\subsection{Bulk and boundary data}

Let $V=\C^2$ be the state space of one site. In tensor-leg notation a subscript
records which factors of a tensor power an operator acts on, so that $R_{12}$
acts on the first two factors of $V^{\otimes3}$ and as the identity on the third.
A bulk matrix of difference form is an invertible
$R(u)\in\operatorname{End}(V\otimes V)$, depending on one complex spectral
parameter, satisfying
\begin{equation}
R_{12}(u-v)R_{13}(u)R_{23}(v)
=R_{23}(v)R_{13}(u)R_{12}(u-v).                 \label{eq:ybe}
\end{equation}
Reading the two sides as amplitudes for three world lines, \cref{eq:ybe} says
that a third line may be slid through the crossing of the other two without
changing the amplitude, which is what makes multi-particle scattering
factorize.

Write $P(x\otimes y)=y\otimes x$ for the exchange of the two factors and
$R_{21}=PR_{12}P$ for the bulk matrix with its two legs interchanged. The bulk
matrix is \emph{regular} when it is holomorphic at the origin with
\begin{equation}
R(0)=P,                                          \label{eq:regularity}
\end{equation}
so that at coincident rapidities the two excitations simply swap places; every
matrix considered below is regular. Regularity makes the logarithmic derivative
of the closed-chain transfer matrix local, and the two-site density it produces
is
\begin{equation}
h=P R'(0).                                      \label{eq:density}
\end{equation}
The matrix $h$ is the interaction of one nearest-neighbour bond: the closed-chain
Hamiltonian is the sum of its copies over the bonds of the ring.

A bulk matrix is $PT$ symmetric when swapping its two spaces has the same effect
as transposing both, $R_{21}(u)=R_{12}^{t_1t_2}(u)$, with $t_1,t_2$ the
transposition on the first or second leg. The six- and eight-vertex matrices are
$PT$ symmetric; most of the six families below are not, and that is why the
ordering of the bulk factors in the next equation has to be stated rather than
assumed.

An integrable reflecting end is a matrix-valued function $K(u)$ on $V$: an
excitation arriving with rapidity $u$ leaves with rapidity $-u$ and its internal
state multiplied by $K(u)$. Throughout this paper the right reflection equation
is the literal identity
\begin{equation}
R_{12}(u-v)K_1(u)R_{21}(u+v)K_2(v)
=K_2(v)R_{12}(u+v)K_1(u)R_{21}(u-v),           \label{eq:re}
\end{equation}
with $K_1=K\otimes\Id$ and $K_2=\Id\otimes K$. Its content is that two
excitations may reflect off the wall in either order, with the appropriate bulk
exchanges in between, and give the same amplitude. Interchanging the roles of
$R_{12}$ and $R_{21}$ in \cref{eq:re} gives a different equation whenever the
bulk matrix is not $PT$ symmetric, and in general a different solution set
\cite{mezincescu1991}; \cref{eq:re} is the equation solved here.

\begin{definition}[Admissible boundary]
\label{def:admissible}
An admissible boundary for a given bulk matrix is a single-valued meromorphic
function $K:\C\to\Mat_2(\C)$, holomorphic at the origin, with
$K(0)=\kappa \Id$ for some $\kappa\ne0$, satisfying \cref{eq:re} as an identity
of meromorphic functions on $\C^2$. No invertibility is required away from the
origin.
\end{definition}

The normalization at the origin is what makes a boundary an endpoint of the same
chain the bulk matrix describes; the meromorphy is global, so a solution defined
only near the origin, or only on a sheet of a branched surface, does not count
until it is shown to continue.

Two admissible boundaries describe the same endpoint when they differ by the
following equivalence. Multiplication by a nonzero scalar meromorphic function
holomorphic and nonzero at the origin is a gauge; it can always be used to
normalize $K(0)=\Id$, and it is used that way from here on. At a fixed bulk
matrix the only further identification is
\begin{equation}
K(u)\sim f(u)\,S\,K(\lambda u)\,S^{-1},                     \label{eq:equivalence}
\end{equation}
where $f$ is such a scalar gauge, $\lambda\in\C^\times$, $S\in GL_2(\C)$, and the
pair $(\lambda,S)$ preserves the same literal bulk matrix projectively:
\begin{equation}
(S\otimes S)\,R(\lambda u)\,(S^{-1}\otimes S^{-1})
=\varphi(u)\,R(u),\qquad \varphi(0)=1,                 \label{eq:bulk-aut}
\end{equation}
for some nonzero meromorphic $\varphi$. Written out, \cref{eq:bulk-aut} says that
rescaling the rapidity and conjugating both legs returns the same bulk matrix up
to an overall function, so that the chain has not been changed. A similarity that
moves the bulk parameters, that is singular on a degeneration, or that depends on
the rapidity is therefore not an equivalence, however natural it may look.

\subsection{The six bulk matrices}
\label{subsec:sixmatrices}

Fix the ordered basis $(e_1\otimes e_1,e_1\otimes e_2,e_2\otimes e_1,e_2\otimes
e_2)$ of $V\otimes V$, in which
\begin{equation}
P=\begin{pmatrix}1&0&0&0\\0&0&1&0\\0&1&0&0\\0&0&0&1\end{pmatrix}.
\label{eq:P}
\end{equation}
This is the exchange operator written out: it moves the second basis vector of
the middle block to the third and back, and fixes the two aligned states.

The six families are the ungraded Classes~1--6 of \cite{deLeeuw2019}, quoted here
from version~3 of that reference. All parameters are complex, and every apparent
parameter singularity is resolved by analytic continuation. Class~1 is subject to
$a_1a_3=a_2a_4$ and reads
\begin{equation}
R_1(u)=\begin{pmatrix}
1&\frac{a_1(e^{a_5u}-1)}{a_5}&\frac{a_2(1-e^{-a_5u})}{a_5}&
\frac{(a_1a_3+a_2a_4)(\cosh(a_5u)-1)}{a_5^2}\\
0&0&e^{-a_5u}&\frac{a_4(1-e^{-a_5u})}{a_5}\\
0&e^{a_5u}&0&\frac{a_3(e^{a_5u}-1)}{a_5}\\
0&0&0&1
\end{pmatrix}.                                      \label{eq:R1}
\end{equation}
Each quotient in \cref{eq:R1} is read at $a_5=0$ as its limit, so that the upper
right entry becomes $(a_1a_3+a_2a_4)u^2/2$ and the exponentials become $1$;
$a_5$ is the rate at which the two off-diagonal amplitudes grow.

For Class~2 put
\begin{equation}
\mathcal N=\begin{pmatrix}
0&a_2&a_3-a_2&a_5\\0&a_1&0&a_4\\
0&0&-a_1&a_3-a_4\\0&0&0&0
\end{pmatrix},
\label{eq:C2-N}
\end{equation}
a nilpotent-on-the-diagonal generator whose square is strictly upper triangular,
and set
\begin{equation}
R_2(u)=uP\left[
\frac{a_1}{\sinh(a_1u)}I_4+\mathcal N+
\frac{\tanh(a_1u/2)}{a_1}\mathcal N^2\right].             \label{eq:R2}
\end{equation}
The bulk matrix of Class~2 is thus the exchange operator dressed by a quadratic
polynomial in one fixed matrix, with rapidity-dependent coefficients; at $a_1=0$
it is $P[I_4+u\mathcal N+u^2\mathcal N^2/2]$. The source prints the leading
scalar without an explicit $I_4$; multiplying by $I_4$ is the dimensionally
consistent reading and the only one compatible with $R_2(0)=P$.

Class~3 is
\begin{equation}
R_3(u)=\begin{pmatrix}
e^{-a_1u}&a_3(e^{(a_1-a_2)u}-e^{-a_1u})&
a_3(e^{(a_1+a_2)u}-e^{-a_1u})&0\\
0&0&e^{(a_1+a_2)u}&0\\
0&e^{(a_1-a_2)u}&0&0\\
0&0&0&e^{-a_1u}
\end{pmatrix},                                      \label{eq:R3}
\end{equation}
a six-vertex-like pattern in which the two middle amplitudes carry independent
exponential rates and $a_3$ couples the aligned state to the two mixed ones.
Class~4 is
\begin{equation}
R_4(u)=\begin{pmatrix}
e^{a_1u}&\frac{a_2\sinh(a_1u)}{a_1}&
\frac{a_2\sinh(a_1u)}{a_1}&
\frac{e^{a_1u}(a_2a_4+a_1a_3\coth(a_1u))\sinh^2(a_1u)}{a_1^2}\\
0&0&e^{-a_1u}&\frac{a_4\sinh(a_1u)}{a_1}\\
0&e^{-a_1u}&0&\frac{a_4\sinh(a_1u)}{a_1}\\
0&0&0&e^{a_1u}
\end{pmatrix};                                      \label{eq:R4}
\end{equation}
its two off-diagonal columns are equal, which is the structural feature that
later forces a conic rather than a line of admissible first derivatives. At
$a_1=0$ its rows are $(1,a_2u,a_2u,a_3u+a_2a_4u^2)$, $(0,0,1,a_4u)$,
$(0,1,0,a_4u)$ and $(0,0,0,1)$.

Class~5 is
\begin{equation}
R_5(u)=(1-a_1u)\begin{pmatrix}
1+2a_1u&a_2u&-a_2u&a_2a_3u^2\\
0&2a_1u&1&-a_3u\\0&1&2a_1u&a_3u\\0&0&0&1+2a_1u
\end{pmatrix},                                      \label{eq:R5}
\end{equation}
and Class~6 is
\begin{equation}
R_6(u)=(1-a_1u)(1+2a_1u)\begin{pmatrix}
1&a_2u&a_2u&-a_2^2u^2(1+2a_1u)\\
0&\frac{2a_1u}{1+2a_1u}&\frac1{1+2a_1u}&-a_2u\\
0&\frac1{1+2a_1u}&\frac{2a_1u}{1+2a_1u}&-a_2u\\
0&0&0&1
\end{pmatrix}.                                      \label{eq:R6}
\end{equation}
These last two are rational rather than exponential: every entry is a polynomial
in the rapidity after the prefactor of \cref{eq:R6} cancels the apparent pole of
its middle block. Each of \cref{eq:R1,eq:R2,eq:R3,eq:R4,eq:R5,eq:R6} satisfies
$R_j(0)=P$, so all six are regular and each defines a closed chain with a
two-site density \cref{eq:density}.

The parameter varieties over which the classification runs are those of
\cref{tab:parameter-varieties}, each carrying the loci at which a family
degenerates.

\begin{table}[t]
\centering
\caption{The bulk parameter varieties, with the loci a complete boundary table
must cover in addition to the generic point.}
\label{tab:parameter-varieties}
\begin{tabularx}{\textwidth}{c X X}
\toprule
Class & Parameter variety & Loci that must be treated separately \\
\midrule
1 & $\mathcal A_1=\{\mathbf a\in\C^5:a_1a_3=a_2a_4\}$ & the analytic $a_5=0$ limit \\
2 & $\mathcal A_2=\C^5$ & the analytic $a_1=0$ limit \\
3 & $\mathcal A_3=\C^3$ & $a_3=0$, and every coincidence of exponential rates \\
4 & $\mathcal A_4=\C^4$ & the analytic $a_1=0$ limit, and the degenerate conic \\
5 & $\mathcal A_5=\C^3$ & $a_2+a_3=0$, and the origin \\
6 & $\mathcal A_6=\C^2$ & $a_2=0$, and the origin \\
\bottomrule
\end{tabularx}
\end{table}

\subsection{The question, and why a list of examples cannot answer it}

For each $j\in\{1,\dots,6\}$ and each $\mathbf a\in\mathcal A_j$, determine every
admissible boundary in the sense of \cref{def:admissible}, together with all
overlaps and degenerations between the resulting families and their orbits under
\cref{eq:equivalence}. The quantifier includes every singular and removable bulk
locus.

That the trivial-looking boundary is already a constraint is seen by substituting
it. With $K(u)=\Id$, \cref{eq:re} becomes
\begin{equation}
R_{12}(u-v)R_{21}(u+v)=R_{12}(u+v)R_{21}(u-v),
\label{eq:scalar-condition}
\end{equation}
an identity between products of bulk matrices at shifted arguments which is
automatic for a $PT$-symmetric bulk matrix and not otherwise. So even the
constant boundary has to be checked class by class; and a list of boundaries that
pass the check, however long, still says nothing about the ones not on it.

At the opposite extreme sits a bulk point at which the equation says almost
nothing, and it is instructive because it is the source of every
arbitrary-function family below.

\begin{example}[The constant fiber]
\label{ex:constant}
If $R(u)=P$ identically, then $R_{21}=P$ as well and \cref{eq:re} collapses to
\begin{equation}
\Id\otimes\bigl(K(u)K(v)-K(v)K(u)\bigr)=0.
\label{eq:constant-fiber}
\end{equation}
The reflection equation has become the statement that the boundary matrix
commutes with itself at any two rapidities, and nothing more.
\end{example}

\begin{lemma}[Centralizer]
\label{lem:centralizer}
Suppose $K:\C\to\Mat_2(\C)$ is meromorphic, that $K(0)=\Id$, and that
$[K(u),K(v)]=0$ holds for all $u$ and $v$. Then either $K$ is scalar, or there is a fixed
nonzero traceless $B_0\in\Mat_2(\C)$ and a meromorphic $\psi$ with $\psi(0)=0$
such that, modulo the scalar gauge,
\begin{equation}
K(u)=\Id+\psi(u)B_0.                  \label{eq:commuting-family}
\end{equation}
Conversely every such $K$ commutes with itself at all rapidities.
\end{lemma}

\begin{proof}
If $K(u_0)$ is nonscalar for some $u_0$, then every $K(v)$ lies in the
centralizer of $K(u_0)$, which for a $2\times2$ matrix is the two-dimensional
span of $\Id$ and $K(u_0)$. Hence $K(v)=\alpha(v)\Id+\beta(v)K(u_0)$ with
meromorphic coefficients, and subtracting the trace and dividing by the scalar
gauge gives \cref{eq:commuting-family} with $B_0$ the traceless part of
$K(u_0)$. Two matrices of that form commute because $B_0$ commutes with itself.
\end{proof}

\Cref{eq:commuting-family} is load-bearing rather than decorative: the free
function $\psi$ is an arbitrary meromorphic function vanishing at the origin, so
a solution may have its first nonzero Taylor coefficient at any order whatever.
No polynomial ansatz, and no truncation of the power series at a fixed order, can
see all of these.

\Cref{tab:symbols} collects the symbols that recur throughout the paper; those
belonging to a single class are declared where that class is treated.

\begin{table}[t]
\centering
\caption{Symbols in force throughout. Class-specific abbreviations are declared
where the class is treated.}
\label{tab:symbols}
\begin{tabularx}{\textwidth}{l X}
\toprule
Symbol & Meaning \\
\midrule
$V$, $P$, $R_j(u)$ & site space $\C^2$; exchange operator \cref{eq:P}; the bulk matrix of class $j$ \\
$\mathcal A_j$, $\mathbf a$ & the bulk parameter variety of class $j$, and a point of it \\
$h$ & the two-site density $PR'(0)$ \\
$K(u)$, $\kappa$ & right boundary, and its value $K(0)=\kappa\Id$ at the origin \\
$k_1,k_2,k_3,k_4$ & the entries of $K(u)$, read row by row \\
$B$, $b,c,d$ & the traceless first derivative of $K$ at the origin, and its entries \cref{eq:B} \\
$f(u)$, $\lambda$, $S$, $\varphi(u)$ & the scalar gauge, spectral scale, similarity and bulk scalar of \cref{eq:equivalence,eq:bulk-aut} \\
$\psi(u)$, $B_0$ & an arbitrary meromorphic function vanishing at $0$, and a fixed traceless direction \\
$E_{ij}$, $\sigma_x,\sigma_y,\sigma_z$ & the matrix units and the Pauli matrices \\
$\mathcal C_{R,B}(u)$ & the $16\times4$ coefficient matrix \cref{eq:derivative-system} \\
$M$, $\rho$, $\chi(u)$ & the crossing matrix, shift and scalar \cref{eq:crossing} \\
$L$, $T_a$, $\widehat T_a$, $t_L(u)$ & chain length, the two monodromies, and the double-row transfer matrix \\
$K_L$, $\mathcal H_L$ & the left boundary and the open-chain Hamiltonian \\
$\Lambda(u)$, $\Gamma(u)$ & a self-commuting family and the entire twist conjugating it \\
\bottomrule
\end{tabularx}
\end{table}

\section{The differentiated system and why it is exhaustive}
\label{sec:exhaustiveness}

\subsection{From the reflection equation to a linear system}

Take a representative normalized by $K(0)=\Id$ and write its entries and its
traceless first derivative as
\begin{equation}
K(u)=\begin{pmatrix}k_1(u)&k_2(u)\\k_3(u)&k_4(u)\end{pmatrix},
\qquad
B:=K'(0)-\tfrac12\tr\bigl(K'(0)\bigr)\Id
=\begin{pmatrix}b&c\\d&-b\end{pmatrix}.       \label{eq:B}
\end{equation}
The four functions $k_1,\dots,k_4$ are unrestricted meromorphic functions: no
entry is set to zero and no diagonal, triangular, polynomial, rational or
exponential form is assumed. The matrix $B$ is the only datum of the boundary
that survives the scalar gauge at first order, and $(b,c,d)$ is called the
\emph{slope} of the boundary. For the derivation the remaining scalar freedom is
used to impose $\tr K(u)=2$ near the origin, which is possible because
$\tr K(0)=2$; then $K'(0)=B$ literally. The displayed representatives in
\cref{sec:classification} sometimes use the more convenient gauge $k_1=1$
instead, and their traceless first derivative is still \cref{eq:B}.

Differentiate \cref{eq:re} with respect to $v$ at $v=0$, using $K(0)=\Id$ and
$K'(0)=B$. Both sides become linear in the single remaining unknown function
$K(u)$, and collecting the coefficient of each entry of $K(u)$ gives a linear
map. Explicitly, for a matrix unit $E_{pq}$,
\begin{equation}
\begin{split}
\Phi_{R,B}(E_{pq})={}&
-2R'(u)(E_{pq}\otimes \Id)R_{21}(u)
+2R(u)(E_{pq}\otimes \Id)R'_{21}(u)\\
&+R(u)(E_{pq}\otimes \Id)R_{21}(u)B_2
-B_2R(u)(E_{pq}\otimes \Id)R_{21}(u),
                                                               \label{eq:derivative-column}
\end{split}
\end{equation}
where $B_2=\Id\otimes B$. Each of the four terms of \cref{eq:derivative-column}
is the contribution of one differentiated factor of \cref{eq:re}: the two bulk
matrices at shifted arguments contribute the derivative terms with opposite
signs, and the boundary at the second leg contributes the commutator with $B$.
Vectorizing $\Phi_{R,B}(E_{pq})\in\Mat_4(\C)$ row by row gives a column of
length sixteen; assembling the four columns in the order
$(E_{11},E_{12},E_{21},E_{22})$ gives a $16\times4$ matrix
$\mathcal C_{R,B}(u)$, and every admissible boundary obeys
\begin{equation}
\mathcal C_{R,B}(u)\,\bigl(k_1(u),k_2(u),k_3(u),k_4(u)\bigr)^{t}=0.
             \label{eq:derivative-system}
\end{equation}
In words: at each rapidity the four entries of an admissible boundary form a
vector in the kernel of an explicitly computable $16\times4$ matrix built from
the bulk data and the boundary slope. \Cref{eq:derivative-system} is necessary,
not sufficient; sufficiency is restored below by substituting each candidate back
into \cref{eq:re}. A scalar factor in $R$ multiplies
\cref{eq:derivative-column} by its square and the two mixed terms cancel, so the
kernel is unchanged by the removal of an overall scalar from the bulk matrix; the
scalars are restored in every direct verification and in every open-chain
computation.

\subsection{The rank trichotomy}

Over the field of meromorphic functions, $\mathcal C_{R,B}(u)$ has a rank, and
that rank is what decides the answer.

\begin{lemma}[Rank trichotomy]
\label{lem:trichotomy}
Fix $j$ and a bulk point $\mathbf a\in\mathcal A_j$, and a slope $(b,c,d)$. Then
$\rank\mathcal C_{R,B}$ over the meromorphic function field is $3$, $2$, or $0$.
\begin{enumerate}
\item If the rank is three, the kernel is one-dimensional. The unique normalized
generator, if it satisfies \cref{eq:re}, is the whole set of admissible
boundaries with that slope, modulo the scalar gauge; if it does not, there is
none.
\item If the rank is two or zero, the equations obtained by specializing
\cref{eq:re} to that cell --- computed there directly, never by dividing through
by a quantity that vanishes on it --- reduce, after conjugation by an entire and
everywhere invertible twist $\Gamma(u)$, to $[\Lambda(u),\Lambda(v)]=0$ for
$K(u)=\Gamma(u)\Lambda(u)\Gamma(u)$. By \cref{lem:centralizer} the admissible
boundaries on that cell are exactly $\Gamma(\Id+\psi B_0)\Gamma$ for a fixed
traceless direction $B_0$ and an arbitrary meromorphic $\psi$ with $\psi(0)=0$,
together with the scalar branch.
\end{enumerate}
\end{lemma}

The lemma is decided cell by cell, and the cells are the content of the next
statement. It is the one step of this paper carried by exhaustive symbolic
computation, so we say exactly what was computed and why it settles the question.

\begin{lemma}[Constructible stratification]
\label{lem:stratification}
For each $j$, the product of the bulk parameter variety $\mathcal A_j$ with the
slope space $\C^3$ is partitioned into finitely many disjoint constructible
cells, listed in \cref{sec:classification}, on each of which
$\rank\mathcal C_{R,B}$ is constant. The cells cover the product exactly; no
residual component remains.
\end{lemma}

\begin{proof}[Proof, by exhaustive symbolic computation]
The object attached to a point of $\mathcal A_j\times\C^3$ is the explicit
$16\times4$ matrix $\mathcal C_{R,B}(u)$ of \cref{eq:derivative-system}, whose
entries are polynomials in the bulk parameters, the slope, and --- depending on
the class --- either the rapidity or the exponentials $e^{\lambda u}$ appearing
in the bulk matrix. The test applied to a point is the rank of that matrix, and
the rank is decided by which of its minors vanish identically in $u$. This is
decisive because \cref{lem:trichotomy} turns each of the three possible ranks
into a complete description of the solution set at that point: at rank three the
kernel has dimension one and there is at most one boundary, at rank two or zero
the equation is a commutator identity whose solutions \cref{lem:centralizer}
exhausts. Nothing about the boundary beyond its slope enters the test.

The cells arise by construction, which is why they cover. Let $I_4$ be the ideal
generated by the sixteen maximal minors of $\mathcal C_{R,B}$, regarded as
polynomials in the bulk parameters and the slope after clearing the rapidity
dependence. Its radical was computed exactly, together with the associated primes
of that radical; the complement of $V(I_4)$ is the rank-three cell, and each
associated prime gives a component on which the rank drops. On each such
component the ideal of $3\times3$ minors was formed anew, from the matrix
specialized to that component rather than from a generic basis, and its radical
computed in turn; this is what distinguishes a genuine lower-rank stratum from an
artefact of dividing by a quantity that vanishes there. The recursion terminates
because each step strictly decreases dimension, and the resulting finitely many
locally closed sets are disjoint after the precedence stated with each table.
Two classes need one further refinement. For Class~3 the minors are exponential
rather than polynomial, so a minor can vanish identically in $u$ through a
coincidence of exponential rates and not through the vanishing of a coefficient;
every such coincidence was solved for, giving the finite list of rate collisions
recorded in \cref{subsec:c3}, and each collision curve was treated as its own
component with the matrix recomputed on it. For Class~4 the rank-drop locus is a
conic in the slope, so it was saturated and covered by both projective charts,
and the two points where the chart used for the generic formula degenerates were
computed separately.

What follows is the classification of \cref{thm:raw}: on each cell of the
resulting partition, one of the two branches of \cref{lem:trichotomy} applies,
and its output is the family printed for that cell.
\end{proof}

Some of the certificates are short enough to print, and they are the ones a
reader is most likely to want to redo. In Class~2 with $a_1\ne0$, writing
$\varepsilon=a_2+a_4-a_3$ and $\omega=a_3\varepsilon-a_1a_5$ and using the
dimensionless slope, the radical of the maximal-minor ideal is
\begin{equation}
\sqrt{I_4}=\langle\, b(b+1)\varepsilon\omega,\ d\varepsilon,\ da_5\,\rangle,
\label{eq:C2-radical}
\end{equation}
whose associated primes are $\langle b,d\rangle$, $\langle b+1,d\rangle$,
$\langle\varepsilon,d\rangle$, $\langle\omega,d\rangle$ and
$\langle\varepsilon,a_5\rangle$. Read as a statement about boundaries,
\cref{eq:C2-radical} says that at $a_1\ne0$ the rank can drop only when the lower
left slope $d$ vanishes together with one of two bulk discriminants or a
distinguished value of $b$, or when the two bulk parameters $\varepsilon$ and
$a_5$ vanish together --- and those five possibilities are exactly the five cells listed for this class in
\cref{subsec:c2}. At $a_1=0$ the same ideal becomes
$\langle ba_3\varepsilon,\ da_3\varepsilon\rangle
=\langle b,d\rangle\cap\langle a_3\rangle\cap\langle\varepsilon\rangle$.
For Class~5 with $a\ne0$ one maximal minor is
\begin{equation}
-2a\,[\,{-u^2pd+u^2rd-2ub-2}\,]
[\,{-2u^2apd+2u^2ard-upd+urd-4a}\,]^2,
\label{eq:C5-minor}
\end{equation}
whose constant term in $u$ is $64a^3$; since that is nonzero whenever $a\ne0$,
the minor cannot vanish identically, so the rank is three for every slope and the
Class-5 boundary set off $a=0$ is a single family. At $a=0$ with $p+r\ne0$ the
constant term of the corresponding row minor is $-8(p+r)^3$, giving the same
conclusion there. For Class~6 with $a\ne0$ the corresponding minor is
\begin{equation}
-64a^3\,(2u^3aQd-ub-1),
\label{eq:C6-minor}
\end{equation}
again with nonzero constant term, so the rank is uniformly three. These four
displays are the reason Classes~5 and~6 have single polynomial families off their
degenerate loci while Classes~1--4 fragment.

\subsection{Global continuation}
\label{subsec:continuation}

The classification is a statement about meromorphic functions on the whole plane,
not about germs, and the step from one to the other is the following.

Each isolated family below is a rational expression in $u$ and in finitely many
exponentials $e^{\lambda u}$, whose denominator is nonzero at $u=0$ on the cell
where it is stated. Such an expression is single-valued and meromorphic on all of
$\C$; the later zeros of its denominator are poles, which \cref{def:admissible}
permits. On a rank-three cell, an arbitrary admissible boundary and the displayed
kernel generator agree on a nonempty open set of rapidities, hence everywhere, by
the identity theorem for meromorphic functions. On a rank-two or rank-zero cell,
the specialized equations determine the fixed traceless direction $B_0$ without
dividing by any coordinate, so the resulting identity persists through the zeros
and poles of the arbitrary function; and the twists $\Gamma$ used there are
entire and everywhere invertible, so conjugating by them neither creates nor
destroys singularities. No branch choice, monodromy, or purely local continuation
is hidden anywhere in the argument.

\section{The classification}
\label{sec:classification}

Throughout this section boundaries are normalized by $K(0)=\Id$ and listed modulo
the scalar gauge, but before the fixed-bulk quotient of \cref{sec:quotient}. In
every isolated family, $(b,c,d)$ is the slope \cref{eq:B}. Every function called
$\psi$, $\psi_1$ or $\psi_2$ is an arbitrary single-valued meromorphic function
on $\C$, holomorphic at the origin, and $\Lambda$ denotes a family commuting with
itself at all rapidities as in \cref{lem:centralizer}.

\begin{theorem}[Classification]
\label{thm:raw}
For every $j\in\{1,\dots,6\}$ and every $\mathbf a\in\mathcal A_j$, the matrices
listed in \cref{subsec:c1,subsec:c2,subsec:c3,subsec:c4,subsec:c5,subsec:c6}, and
no others, are the admissible boundaries of \cref{def:admissible} for $R_j$ at
$\mathbf a$. The cells are finite in number, constructible, disjoint after the
stated precedence, and cover the full parameter variety. Each family of arbitrary
functions contains its own finite specializations, and the overlap statements say
where two families meet.
\end{theorem}

The proof is given after the six subsections, once the objects it refers to have
been displayed.

\subsection{Class 1}
\label{subsec:c1}

Assume $a_1a_3=a_2a_4$, write $\theta=e^{a_5u}$ for the bulk exponential, and set
\begin{equation}
\mathcal E_1=\{a_1=a_4,\ a_2=a_3\}.
\label{eq:C1-E1}
\end{equation}
The locus $\mathcal E_1$ is where the two off-diagonal columns of $R_1$ acquire
matching rates, and it is the only Class-1 locus carrying arbitrary functions at
$a_5\ne0$. For $a_5\ne0$ define
\begin{equation}
K_+(u)=\begin{pmatrix}
1+\frac b{a_5}(1-\theta^{-1})&\frac c{2a_5}(\theta-\theta^{-1})\\
0&1-\frac b{a_5}(\theta-1)
\end{pmatrix},                                      \label{eq:C1-plus}
\end{equation}
\begin{equation}
K_-(u)=\begin{pmatrix}
\theta^{-2}\frac{(b+2a_5)\theta-b-a_5}{a_5}&0\\
\frac d{2a_5}(\theta-\theta^{-1})&
\theta\frac{-(b+a_5)\theta+b+2a_5}{a_5}
\end{pmatrix}.                                      \label{eq:C1-minus}
\end{equation}
These are the upper- and lower-triangular boundaries whose off-diagonal entry
grows like the bulk amplitude it must follow; each is determined by its slope
alone. On $\mathcal E_1$ let
\begin{equation}
\Gamma_1(u)=\theta^{-1/2}
\begin{pmatrix}1&-\frac{a_2}{a_5}(\theta-1)\\0&\theta\end{pmatrix},
\label{eq:C1-twist}
\end{equation}
an entire and everywhere invertible matrix function, since $\theta^{-1/2}
=e^{-a_5u/2}$ and the determinant is $1$. The complete family there is
\begin{equation}
K(u)=\Gamma_1(u)\Lambda(u)\Gamma_1(u),\qquad [\Lambda(u),\Lambda(v)]=0,
\label{eq:C1-factor-nz}
\end{equation}
so on $\mathcal E_1$ the reflection equation says only that a twisted version of
the boundary commutes with itself, and every self-commuting family is allowed.
The traceless slope of $\Lambda$ has coordinates $(b+a_5,\,c+2a_2,\,d)$.

Off $\mathcal E_1$ the remaining cells are those of \cref{tab:C1-nonzero}.

\begin{table}[t]
\centering\small
\caption{Class~1 at $a_5\ne0$, off the locus $\mathcal E_1$.}
\label{tab:C1-nonzero}
\begin{tabularx}{\textwidth}{>{\RaggedRight\arraybackslash}X >{\RaggedRight\arraybackslash}X}
\toprule
Bulk and slope cell & Complete family \\
\midrule
$(a_2,a_3)\ne(0,0)$ and $(a_1,a_4)\ne(0,0)$ & $d=0$ forced; uniquely $K_+$ \\
$(a_2,a_3)\ne(0,0)$, $(a_1,a_4)=(0,0)$, $b\ne-a_5$ & $d=0$ forced; uniquely $K_+$ \\
the same bulk cell with $b=-a_5$ &
$\psi_1(u)\diag(1,e^{2a_5u})+\psi_2(u)E_{12}$, with
$\psi_1(0)=1$, $\psi_2(0)=0$, $\psi_1'(0)=-a_5$, $\psi_2'(0)=c$ \\
$(a_2,a_3)=(0,0)$, $d=0$, $c\ne0$ & uniquely $K_+$ \\
$(a_2,a_3)=(0,0)$, $c=0$, $d\ne0$ & uniquely $K_-$ \\
$(a_2,a_3)=(0,0)$, $c=d=0$ & every normalized diagonal family, with traceless slope $\diag(b,-b)$ \\
\bottomrule
\end{tabularx}
\end{table}

For $a_5=0$ the exponentials degenerate and the boundary becomes polynomial:
\begin{equation}
K_0(u)=\begin{pmatrix}
1+bu-\frac{a_2+a_3}{2}du^2&cu+a_2a_3du^3\\
du&1-bu-\frac{a_2+a_3}{2}du^2
\end{pmatrix},                                      \label{eq:C1-zero}
\end{equation}
a cubic in the rapidity whose higher coefficients are fixed by the slope and the
bulk parameters. Every slope occurs. Off $\mathcal E_1$, \cref{eq:C1-zero} is the
unique family except in two cases: at $(a_1,a_4)=(0,0)$ with $(a_2,a_3)\ne(0,0)$
and $b=d=0$, every $\psi_1\Id+\psi_2E_{12}$ with $\psi_1(0)=1$, $\psi_2(0)=0$,
$\psi_1'(0)=0$, $\psi_2'(0)=c$ occurs; and at $(a_2,a_3)=(0,0)$ with
$(a_1,a_4)\ne(0,0)$ and $c=d=0$, every normalized diagonal family occurs. On
$\mathcal E_1$, with $\Gamma_1(u)=\Id-a_2uE_{12}$, the complete family is again
\begin{equation}
K(u)=\Gamma_1(u)\Lambda(u)\Gamma_1(u),\qquad [\Lambda(u),\Lambda(v)]=0,
\label{eq:C1-factor-zero}
\end{equation}
the $a_5\to0$ limit of \cref{eq:C1-factor-nz}, with auxiliary slope
$(b,c+2a_2,d)$. The finite formulas above are contained in these
arbitrary-function families wherever the two overlap, so no boundary is counted
twice.

\subsection{Class 2}
\label{subsec:c2}

Set
\begin{equation}
\varepsilon=a_2+a_4-a_3,\qquad \omega=a_3\varepsilon-a_1a_5,
\label{eq:C2-vars}
\end{equation}
the two discriminants that \cref{eq:C2-radical} showed to control the rank drops:
$\varepsilon$ measures the failure of the three linear parameters to balance, and
$\omega$ couples that failure to the quadratic parameter $a_5$. For $a_1\ne0$
define
\begin{equation}
K_\omega(u)=\begin{pmatrix}
1+\frac b{a_1}(1-e^{-a_1u})&\frac c{a_1}\sinh(a_1u)\\
0&1-\frac b{a_1}(e^{a_1u}-1)
\end{pmatrix},                                      \label{eq:C2-omega}
\end{equation}
\begin{equation}
K_\varepsilon(u)=\begin{pmatrix}
\cosh(a_1u)+\frac b{a_1}\sinh(a_1u)&
\frac{\sinh(a_1u)}{a_1^2}\bigl\{(a_1c-2a_4b)\cosh(a_1u)+2a_4b\bigr\}\\
0&\cosh(a_1u)-\frac b{a_1}\sinh(a_1u)
\end{pmatrix}.                                      \label{eq:C2-delta}
\end{equation}
Both are upper triangular with unit determinant to leading order, and they differ
in which discriminant vanishes: $K_\omega$ is the boundary on $\omega=0$ and
$K_\varepsilon$ the boundary on $\varepsilon=0$. The nonzero-rate partition is
\cref{tab:C2-nonzero}.

\begin{table}[t]
\centering\small
\caption{Class~2 at $a_1\ne0$.}
\label{tab:C2-nonzero}
\begin{tabularx}{\textwidth}{>{\RaggedRight\arraybackslash}X >{\RaggedRight\arraybackslash}X}
\toprule
Bulk and slope cell & Complete family \\
\midrule
$\varepsilon=0$, $a_5=0$ & the twisted commuting family \cref{eq:C2-factor}; every slope \\
$\varepsilon=0$, $a_5\ne0$, $b\ne-a_1$ & $d=0$ forced; uniquely $K_\varepsilon$ \\
$\varepsilon=0$, $a_5\ne0$, $b=-a_1$ &
$\begin{psmallmatrix}e^{-a_1u}&\psi(u)\\0&e^{a_1u}\end{psmallmatrix}$,
$\psi(0)=0$, $\psi'(0)=c$ \\
$\varepsilon\ne0$, $\omega=0$ & $d=0$ forced; uniquely $K_\omega$, for arbitrary $b,c$ \\
$\varepsilon\omega\ne0$ & $d=0$ and $b\in\{0,-a_1\}$ forced; uniquely $K_\omega$ \\
\bottomrule
\end{tabularx}
\end{table}

On the first row of \cref{tab:C2-nonzero} put
\begin{equation}
\Xi=\begin{pmatrix}0&a_4\\0&-a_1\end{pmatrix},\qquad
\Gamma_2(u)=e^{-a_1u/2}\,e^{-u\Xi},
\label{eq:C2-twist}
\end{equation}
again entire and everywhere invertible, and then
\begin{equation}
K(u)=\Gamma_2(u)\Lambda(u)\Gamma_2(u),\qquad [\Lambda(u),\Lambda(v)]=0,
\label{eq:C2-factor}
\end{equation}
with auxiliary slope $(b+a_1,\,c+2a_4,\,d)$: on that cell the bulk matrix
factorizes through the exchange operator and an entire twist, so the reflection
equation again reduces to self-commutation.

For $a_1=0$ define the two polynomial boundaries
\begin{equation}
K_{\varepsilon,0}(u)=\begin{pmatrix}
1+bu-a_4du^2&cu+a_4^2du^3\\du&1-bu-a_4du^2
\end{pmatrix},                                      \label{eq:C2-delta0}
\end{equation}
\begin{equation}
K_{a_3,0}(u)=\begin{pmatrix}
1+bu+\frac{a_2-a_4}{2}du^2&cu-a_2a_4du^3\\
du&1-bu+\frac{a_2-a_4}{2}du^2
\end{pmatrix},                                      \label{eq:C2-a30}
\end{equation}
whose quadratic and cubic coefficients differ exactly by which bulk combination
survives the limit. The zero-rate partition is \cref{tab:C2-zero}.

\begin{table}[t]
\centering\small
\caption{Class~2 at $a_1=0$.}
\label{tab:C2-zero}
\begin{tabularx}{\textwidth}{>{\RaggedRight\arraybackslash}X >{\RaggedRight\arraybackslash}X}
\toprule
Cell & Complete family \\
\midrule
$\varepsilon=0$, $a_5=0$ &
$(\Id-a_4uE_{12})\Lambda(u)(\Id-a_4uE_{12})$ with $[\Lambda(u),\Lambda(v)]=0$; auxiliary slope $(b,c+2a_4,d)$ \\
$\varepsilon=0$, $a_5\ne0$, $(b,d)\ne(0,0)$ & uniquely $K_{\varepsilon,0}$; $c$ arbitrary \\
$\varepsilon=0$, $a_5\ne0$, $b=d=0$ & $\Id+\psi(u)E_{12}$, $\psi(0)=0$, $\psi'(0)=c$ \\
$a_3=0$, $\varepsilon\ne0$ & uniquely $K_{a_3,0}$; every slope \\
$a_3\varepsilon\ne0$ & uniquely $\Id+cuE_{12}$ \\
\bottomrule
\end{tabularx}
\end{table}

\subsection{Class 3}
\label{subsec:c3}

Class~3 is the class in which exponential rates rather than polynomial
coefficients control the answer. Put
\begin{equation}
m=2a_1-a_2,\qquad n=2a_1+a_2,\qquad g=a_3,
\qquad X=e^{mu},\quad Y=e^{nu},                    \label{eq:C3-vars}
\end{equation}
so that $m$ and $n$ are the two independent rates carried by the middle block and
$g$ is the coupling of the aligned state to the mixed ones. Removing the harmless
overall scalar $e^{-a_1u}$ leaves
\begin{equation}
\widehat R_3(u)=\begin{pmatrix}
1&g(X-1)&g(Y-1)&0\\0&0&Y&0\\0&X&0&0\\0&0&0&1
\end{pmatrix},                                      \label{eq:C3-bulk}
\end{equation}
which shows that the whole class depends only on $(m,n,g)$, and that the bulk
matrix is the exchange operator exactly when $m=n=0$. The complete partition is
\cref{tab:C3-cells}, the $g\ne0$ cells taking precedence over the $g=0$ ones.

\begin{table}[t]
\centering\footnotesize
\caption{The complete constructible partition for Class~3.}
\label{tab:C3-cells}
\begin{tabularx}{\textwidth}{>{\RaggedRight\arraybackslash}X >{\RaggedRight\arraybackslash}X}
\toprule
Condition & Types of family that occur \\
\midrule
$g\ne0$, $mn(m-n)(m+n)(n+2m)(m+2n)\ne0$ & diagonal, and mixed upper \\
$g\ne0$, $m=n\ne0$ & diagonal, lower, pure upper, mixed upper \\
$g\ne0$, $n=-m\ne0$ & one full upper-triangular component \\
$g\ne0$, $n=-2m\ne0$ & the inherited diagonal and mixed branches, plus one extra component \\
$g\ne0$, $m=-2n\ne0$ & the inherited diagonal and mixed branches, plus one extra component \\
$g\ne0$, $m=0$, $n\ne0$ & unique, off one arbitrary-function stabilizer line \\
$g\ne0$, $n=0$, $m\ne0$ & unique, off an arbitrary diagonal stabilizer \\
$m=n=0$, any $g$ & every self-commuting family \\
$g=0$, $m+n\ne0$ & unique, off an arbitrary diagonal stabilizer \\
$g=0$, $m=-n\ne0$ & arbitrary diagonal, plus a homogeneous projective family \\
\bottomrule
\end{tabularx}
\end{table}

On the generic cell the two families are
\begin{equation}
K_D(u)=\diag\left(1,
\frac{(n-b)Y+b}{Y\{b(Y-1)+n\}}\right)            \label{eq:C3-diag}
\end{equation}
and, for $d=0$ and $b=-c/(2g)$,
\begin{equation}
K_U(u)=\begin{pmatrix}
1&\dfrac{cg(X^2-1)}{2mgX-c(X-1)}\\[2mm]
0&\dfrac{X\{2mg+c(X-1)\}}{2mgX-c(X-1)}
\end{pmatrix},                                      \label{eq:C3-mixed}
\end{equation}
with $b$ free in \cref{eq:C3-diag} and $c$ free in \cref{eq:C3-mixed}. The
diagonal branch uses only the rate $n$ and is available for every $b$; the mixed
branch exists only when the slope satisfies the one linear relation $2gb+c=0$,
which ties the upper entry of the boundary to the bulk coupling.

On $m=n=\mu\ne0$ with $\theta=e^{\mu u}$, the diagonal and mixed branches are the
specializations of \cref{eq:C3-diag,eq:C3-mixed}, and two further branches open:
\begin{equation}
K_{\mathrm{lo}}(u)=\begin{pmatrix}
1&0\\[1mm]
\dfrac{d(\theta^2-1)}{gd(\theta-1)+2\mu\theta}&
\dfrac{\theta\{2\mu-gd(\theta-1)\}}{gd(\theta-1)+2\mu\theta}
\end{pmatrix},\qquad c=0,\quad b=\frac{gd}{2},       \label{eq:C3-lower}
\end{equation}
\begin{equation}
K_{PU}(u)=\Id+\frac{c(\theta^2-1)}{2\mu\theta}E_{12},\qquad b=d=0.
\label{eq:C3-pure-upper}
\end{equation}
When the two rates coincide the class regains the symmetry between the upper and
the lower triangle, which is why a lower-triangular boundary exists there and
nowhere else at $g\ne0$.

The three cells on which the rates are in negative ratio are not optional, and
they are the reason the exponential grouping of \cref{lem:stratification} is
needed. With $\theta=e^{mu}$ and $n=-m\ne0$, every slope with $d=0$ and arbitrary
$b,c$ has the unique solution
\begin{equation}
K_{-1}(u)=\begin{pmatrix}
1&\dfrac{c(\theta^2-1)}{2\{(m+b)\theta-b\}}\\[2mm]
0&\dfrac{\theta\{m+b-b\theta\}}{(m+b)\theta-b}
\end{pmatrix}.                                      \label{eq:C3-neg1}
\end{equation}
If $n=-2m\ne0$ the extra component is
\begin{equation}
d=0,\quad b=-m,\qquad
K_{-2}(u)=\begin{pmatrix}1&c(\theta^2-1)/(2m)\\0&\theta^2\end{pmatrix},
\quad \theta=e^{mu},                                     \label{eq:C3-neg2}
\end{equation}
and if $m=-2n\ne0$ it is
\begin{equation}
d=0,\quad b=n,\qquad
K_{-1/2}(u)=\begin{pmatrix}1&c(\theta^2-1)/(2n\theta^2)\\0&\theta^{-2}\end{pmatrix},
\quad \theta=e^{nu}.                                     \label{eq:C3-neghalf}
\end{equation}
Each of \cref{eq:C3-neg1,eq:C3-neg2,eq:C3-neghalf} is a boundary that exists only
on its own curve in the rate plane and is invisible from the generic formulas.

\begin{remark}
\label{rem:C3-resonance}
The necessity of these cells is not a formality. At $g=m=1$ and $n=-1$ the matrix
$\Id+\sinh(u)E_{12}$ is an admissible boundary, and it is a specialization of
\cref{eq:C3-neg1} and of no generic-rate formula. A table built by solving the
system once at generic rates and then substituting would omit it. Solving for the
identical vanishing of the exponential minors shows that this happens for
$n/m\in\{-2,-1,-\tfrac12\}$ and for no other ratio, which is what makes the list
of three complete.
\end{remark}

On $g\ne0$, $m=0$, $n=\mu\ne0$, put $\theta=e^{\mu u}$ and
\begin{equation}
\mathcal D_3=2\mu\theta+2b\theta(\theta-1)-gd(\theta-1)^2 .
\label{eq:C3-D3}
\end{equation}
Off the line $d=0$, $b=-c/(2g)$ the unique boundary is
\begin{equation}
K_{m=0}(u)=\begin{pmatrix}
1&\dfrac{c(\theta^2-1)}{\mathcal D_3}\\[1mm]
\dfrac{d(\theta^2-1)}{\mathcal D_3}&
\dfrac{2\mu\theta-2b(\theta-1)-gd(\theta-1)^2}{\mathcal D_3}
\end{pmatrix},                                      \label{eq:C3-mzero}
\end{equation}
the first boundary in this class with both off-diagonal entries nonzero, and on
the excluded line the family is
\begin{equation}
K(u)\sim \Id+\psi(u)\begin{pmatrix}0&1\\0&1/g\end{pmatrix},
\qquad \psi(0)=0,\quad \psi'(0)=c.                        \label{eq:C3-mzero-fun}
\end{equation}
On $g\ne0$, $n=0$, $m=\mu\ne0$, put $\theta=e^{\mu u}$ and
$\mathcal D_4=b(\theta-1)+\mu\theta$; off $c=d=0$,
\begin{equation}
K_{n=0}(u)=\begin{pmatrix}
1&\dfrac{c(\theta^2-1)}{2\mathcal D_4}\\[1mm]
\dfrac{d(\theta^2-1)}{2\mathcal D_4}&
\dfrac{\theta\{\mu-b(\theta-1)\}}{\mathcal D_4}
\end{pmatrix},                                     \label{eq:C3-nzero}
\end{equation}
while on $c=d=0$ every $\diag(1,\psi(u))$ with $\psi(0)=1$ and $\psi'(0)=-2b$
occurs. When $m=n=0$, \cref{eq:C3-bulk} is the exchange operator for every $g$,
and \cref{ex:constant,lem:centralizer} give the answer.

Two cells remain, both at $g=0$. If $m+n\ne0$, set $Z=XY=e^{(m+n)u}$ and
$\mathcal L=(b+m)Z+n-b$. Off $c=d=0$ the unique boundary is
\begin{equation}
K_{g=0}(u)=\begin{pmatrix}
1&\dfrac{c(Z^2-1)}{2Y\mathcal L}\\[1mm]
\dfrac{d(Z^2-1)}{2Y\mathcal L}&
\dfrac{X\{(n-b)Z+b+m\}}{Y\mathcal L}
\end{pmatrix},                                     \label{eq:C3-gzero}
\end{equation}
which depends on the two rates only through their sum in the numerators and
through $X/Y$ in the corner; on $c=d=0$ every $\diag(1,\psi(u))$ with $\psi(0)=1$
and $\psi'(0)=-2b$ occurs. If instead $m=-n=\mu\ne0$, put $\theta=e^{\mu u}$;
besides the diagonal families the complete nondiagonal family is homogeneous,
\begin{equation}
K(u)\sim\begin{pmatrix}
1&\gamma_1\psi(u)\\ \gamma_2\psi(u)&\theta^2+\nu\,\theta\,\psi(u)
\end{pmatrix},\qquad [\gamma_1:\gamma_2]\in\mathbb P^1,\quad \psi(0)=0,
\label{eq:C3-opposite}
\end{equation}
a one-function family whose off-diagonal direction is a free point of the
projective line rather than a fixed matrix. Its slope satisfies
$c=\gamma_1\psi'(0)$, $d=\gamma_2\psi'(0)$ and $b=-\mu-\nu\psi'(0)/2$, and the
presentation carries the redundancy
$(\gamma_1,\gamma_2,\nu,\psi)\sim(\varsigma\gamma_1,\varsigma\gamma_2,\varsigma\nu,
\psi/\varsigma)$ for $\varsigma\ne0$. The chart $\gamma_1=0$, $\gamma_2\ne0$ is
included explicitly, and no division by either off-diagonal coordinate is used,
so the purely lower-triangular members are not lost.

\subsection{Class 4}
\label{subsec:c4}

For $a_1\ne0$ introduce dimensionless bulk variables and a dimensionless slope,
\begin{equation}
A=\frac{a_2}{a_1},\quad G=\frac{a_3}{a_1},\quad
H=\frac{a_4}{a_1},\quad \theta=e^{a_1u},\qquad
\frac B{a_1}=\begin{pmatrix}b_0&c_0\\d_0&-b_0\end{pmatrix},
\label{eq:C4-vars}
\end{equation}
and the two discriminants
\begin{equation}
\Sigma=A+H,\qquad J=4G-(H-A)^2.                         \label{eq:C4-discriminants}
\end{equation}
The first measures the imbalance of the two off-diagonal amplitudes and the
second is the discriminant of the conic that the admissible slopes lie on. The
nine cells are those of \cref{tab:C4-cells}.

\begin{table}[t]
\centering\footnotesize
\caption{The complete constructible partition for Class~4.}
\label{tab:C4-cells}
\begin{tabularx}{\textwidth}{>{\RaggedRight\arraybackslash}X >{\RaggedRight\arraybackslash}X}
\toprule
Bulk condition & Types of family that occur \\
\midrule
$a_1\ne0$, $\Sigma\ne0$, $J\ne0$ & upper, mixed, the generic conic, and two vertical charts \\
$a_1\ne0$, $\Sigma\ne0$, $J=0$ & upper, mixed, and one genuine degenerate line \\
$a_1\ne0$, $H=-A$, $G\ne A^2$ & a single formula valid for every slope \\
$a_1\ne0$, $H=-A$, $G=A^2$, $A\ne0$ & that formula, plus an arbitrary-function stabilizer \\
$a_1\ne0$, $A=G=H=0$ & that formula, plus an arbitrary diagonal family \\
$a_1=0$, $a_2\ne a_4$ & a single formula valid for every slope \\
$a_1=0$, $a_2=a_4$, $a_3\ne0$ & that formula, plus an arbitrary upper family \\
$a_1=0$, $a_2=a_4\ne0$, $a_3=0$ & that formula, and twisted self-commuting families \\
$a_1=a_2=a_3=a_4=0$ & every self-commuting family \\
\bottomrule
\end{tabularx}
\end{table}

On $a_1\ne0$, $\Sigma\ne0$, the upper and mixed templates are
\begin{equation}
K_U^{(4)}=\Id+\frac{c_0(\theta^2-\theta^{-2})}{4}E_{12},\qquad
b_0=d_0=0,                                    \label{eq:C4-upper}
\end{equation}
and, with $4b_0=(H-A)d_0$, $4c_0=-Gd_0$ and
$\mathcal D_m=d_0\{A\theta^4-\Sigma\theta^2+H\}-8\theta^2$,
\begin{equation}
K_M^{(4)}=\begin{pmatrix}
1&\dfrac{Gd_0(\theta^4-1)}{2\mathcal D_m}\\[2mm]
-\dfrac{2d_0(\theta^4-1)}{\mathcal D_m}&
\dfrac{d_0\{H\theta^4-\Sigma\theta^2+A\}-8\theta^2}{\mathcal D_m}
\end{pmatrix}.                                      \label{eq:C4-mixed}
\end{equation}
\Cref{eq:C4-upper} is the boundary with no diagonal deformation at all;
\cref{eq:C4-mixed} is the one whose three slope coordinates are locked to a
single free parameter by the two displayed relations.

When $J\ne0$ the remaining slopes satisfy
\begin{equation}
d_0=0,\qquad
\mathcal Q:=Gb_0^2+(H-A)b_0c_0+c_0^2-4(AH+G)=0,   \label{eq:C4-conic}
\end{equation}
so the admissible first derivatives form a conic in the $(b_0,c_0)$ plane rather
than a line: this is the structural consequence of the two equal off-diagonal
columns of \cref{eq:R4}. Away from the two points
$(b_0,c_0)=(2,-2H)$ and $(-2,-2A)$, set
\begin{align}
\mathcal E={}&Hb_0^2\theta^4+b_0c_0\theta^4+2Hb_0\theta^4+2c_0\theta^4
-\Sigma b_0^2\theta^2+Ab_0^2-b_0c_0+4\Sigma\theta^2-2Ab_0+2c_0,\nonumber\\
\mathcal N_y={}&-AHb_0^2\theta^4+(H-A)b_0c_0\theta^4+c_0^2\theta^4
+2AHb_0^2\theta^2+(A-H)b_0c_0\theta^2+2\Sigma c_0\theta^2\nonumber\\
&-AHb_0^2+(H-A)b_0c_0+c_0^2,\nonumber\\
\mathcal N_w={}&-Ab_0^2\theta^4+b_0c_0\theta^4+2Ab_0\theta^4-2c_0\theta^4
+\Sigma b_0^2\theta^2-Hb_0^2-b_0c_0-4\Sigma\theta^2-2Hb_0-2c_0,
\label{eq:C4-conic-aux}
\end{align}
three polynomials in $\theta$ whose coefficients are the conic data; then
\begin{equation}
K_C^{(4)}=\begin{pmatrix}
1&\dfrac{(\theta^4-1)\mathcal N_y}{2\theta^2\mathcal E}\\[2mm]
0&-\dfrac{\mathcal N_w}{\mathcal E}
\end{pmatrix}.                                      \label{eq:C4-conic-template}
\end{equation}
The two excluded points are not missing solutions but points at which this chart
degenerates; they are covered by writing, with $\Theta=AH+G$,
\begin{equation}
\Pi_+=4\Theta\theta^4-2\Sigma^2\theta^2+\Sigma^2,\qquad
\Pi_-=\Sigma^2\theta^4-2\Sigma^2\theta^2+4\Theta,
\label{eq:C4-Pi}
\end{equation}
\begin{align}
\Upsilon_+&=2\Sigma\Theta(\theta^4+1)-(5A^2H+4AG+2AH^2+H^3)\theta^2,\nonumber\\
\Upsilon_-&=2\Sigma\Theta(\theta^4+1)-(A^3+2A^2H+5AH^2+4GH)\theta^2,
\label{eq:C4-Upsilon}
\end{align}
so that the two boundaries there are
\begin{equation}
K_{C,+}^{(4)}=\begin{pmatrix}1&-\dfrac{(\theta^4-1)\Upsilon_+}{2\theta^2\Pi_+}\\
0&\Pi_-/\Pi_+\end{pmatrix},
\quad
K_{C,-}^{(4)}=\begin{pmatrix}1&-\dfrac{(\theta^4-1)\Upsilon_-}{2\theta^2\Pi_-}\\
0&\Pi_+/\Pi_-\end{pmatrix}.
                                                               \label{eq:C4-vertical}
\end{equation}
Their diagonal entries are reciprocal, which is the signature of the two vertical
directions of the conic.

When $J=0$ the conic degenerates into two lines, and exactly one of them is
consistent with the actual first derivative:
\begin{equation}
d_0=0,\qquad c_0=-\frac{H-A}{2}b_0+\Sigma.       \label{eq:C4-line-condition}
\end{equation}
Writing
\begin{align}
\mathcal D_\ell&=b_0(\theta^2-1)+2(\theta^2+1),&
\mathcal N_\ell&=b_0(\theta^2-1)-2(\theta^2+1),\nonumber\\
\mathcal F_\ell&=\Sigma b_0^2(\theta^2-1)^2-4\Sigma(\theta^2+1)^2
+8(H-A)b_0\theta^2,
\label{eq:C4-line-aux}
\end{align}
the boundary on that line is
\begin{equation}
K_\ell^{(4)}=\begin{pmatrix}
1&-\dfrac{(\theta^4-1)\mathcal F_\ell}{4\theta^2\mathcal D_\ell^2}\\[2mm]
0&\dfrac{\mathcal N_\ell^2}{\mathcal D_\ell^2}
\end{pmatrix}.                                      \label{eq:C4-line}
\end{equation}
The second algebraic line carries the wrong first derivative except at the
already listed upper point, so it contributes no family; this is a case in which
the derivative condition is necessary but not sufficient, and the check is the
substitution back into \cref{eq:re}.

On $\Sigma=0$ every slope is admissible and one rational template covers them all:
\begin{equation}
K_{\mathrm{ex}}^{(4)}=\begin{pmatrix}
1&(\theta^4-1)\mathcal Y_e/\mathcal D_e\\
-2d_0(\theta^8-1)/\mathcal D_e&-\mathcal W_e/\mathcal D_e
\end{pmatrix},                                      \label{eq:C4-ex}
\end{equation}
with
\begin{align}
\mathcal D_e={}&Ad_0\theta^8-2Ad_0\theta^6-4b_0\theta^6-8\theta^6
+2Ad_0\theta^2+4b_0\theta^2-Ad_0-8\theta^2,\nonumber\\
\mathcal Y_e={}&A^2d_0\theta^4+2Ab_0\theta^4-2A^2d_0\theta^2-2c_0\theta^4
-4Ab_0\theta^2+A^2d_0+2Ab_0-2c_0,\nonumber\\
\mathcal W_e={}&Ad_0\theta^8-2Ad_0\theta^6-4b_0\theta^6+8\theta^6
+2Ad_0\theta^2+4b_0\theta^2-Ad_0+8\theta^2.
\label{eq:C4-ex-aux}
\end{align}
Since $\mathcal D_e=-16$ at $\theta=1$, the chart is regular at the origin, which
is what makes \cref{eq:C4-ex} a legitimate normalized representative. On
$H=-A$, $G=A^2$, $A\ne0$ the stabilizer locus $d_0=0$, $c_0=Ab_0$ enlarges to
\begin{equation}
K(u)\sim \Id+\psi(u)\begin{pmatrix}0&1\\0&-2/A\end{pmatrix},
\qquad \psi(0)=0,\quad \psi'(0)=a_1c_0,                \label{eq:C4-ex-fun}
\end{equation}
and at $A=G=H=0$ the corresponding enlargement is every $\diag(1,\psi(u))$ with
$\psi(0)=1$ and $\psi'(0)=-2b$.

For the removable limit $a_1=0$, write $A=a_2$, $G=a_3$, $H=a_4$ and
$\mathcal D_0=(A+H)du^2-2bu-2$. Away from the functional enlargements, every
slope has the single rational boundary
\begin{equation}
K_0^{(4)}(u)=\begin{pmatrix}
1&-\dfrac{2u(AHdu^2+c)}{\mathcal D_0}\\[2mm]
-\dfrac{2du}{\mathcal D_0}&
\dfrac{(A+H)du^2+2bu-2}{\mathcal D_0}
\end{pmatrix},                                      \label{eq:C4-zero}
\end{equation}
a ratio of quadratics rather than a polynomial, because the limit of the
exponential chart keeps its denominator. If $A=H$ and $G\ne0$, the locus $b=d=0$
enlarges to $\Id+\psi(u)E_{12}$ with $\psi(0)=0$ and $\psi'(0)=c$. If $A=H$ and
$G=0$, put $\Gamma_4(u)=\Id+AuE_{12}$; the complete family is
\begin{equation}
K(u)=\Gamma_4(-u)\{\psi_1(u)\Id+\psi_2(u)B_0\}\Gamma_4(-u),
\label{eq:C4-zero-commuting}
\end{equation}
for a fixed nonscalar $B_0$, together with the scalar branch, where $\psi_1(0)=1$
and $\psi_2(0)=0$; writing
$B_0=\begin{psmallmatrix}B_{11}&B_{12}\\B_{21}&-B_{11}\end{psmallmatrix}$, its
slope is $b=\psi_2'(0)B_{11}$, $c=\psi_2'(0)B_{12}-2A$,
$d=\psi_2'(0)B_{21}$. At $A=G=H=0$ this is \cref{ex:constant} again. The
single-formula branch \cref{eq:C4-zero} meets \cref{eq:C4-zero-commuting}
exactly at
\begin{equation}
B_0=\begin{pmatrix}b&c+2A\\d&-b\end{pmatrix},\qquad
\psi_1(u)=\frac1{1+bu-Adu^2},\qquad \psi_2(u)=u\,\psi_1(u),
\label{eq:C4-overlap}
\end{equation}
so the two descriptions agree on their intersection and nothing is counted twice.

\subsection{Class 5}
\label{subsec:c5}

Write $a=a_1$, $p=a_2$ and $r=a_3$. On the three cells where the rank is three,
one cubic polynomial exhausts the boundaries:
\begin{equation}
K_5(u)=\begin{pmatrix}
1+bu+\frac{p-r}{2}du^2&cu-pr\,du^3\\
du&1-bu+\frac{p-r}{2}du^2
\end{pmatrix}.                                      \label{eq:C5-poly}
\end{equation}
Every entry of \cref{eq:C5-poly} is determined by the slope and the bulk
parameters, and every slope occurs, so the Class-5 boundary set off the
degenerate loci is a three-parameter family of polynomial matrices. The complete
partition is \cref{tab:C5-cells}.

\begin{table}[t]
\centering\small
\caption{The complete Class-5 partition.}
\label{tab:C5-cells}
\begin{tabularx}{\textwidth}{>{\RaggedRight\arraybackslash}X >{\RaggedRight\arraybackslash}X}
\toprule
Condition & Complete family \\
\midrule
$a\ne0$, $(p,r)\ne(0,0)$ & $K_5$, every slope \\
$a\ne0$, $p=r=0$ & $K_5=\Id+uB$, every slope \\
$a=0$, $p+r\ne0$ & $K_5$, every slope \\
$a=0$, $r=-p\ne0$ & twisted self-commuting families \cref{eq:C5-twist} \\
$a=p=r=0$ & every self-commuting family \\
\bottomrule
\end{tabularx}
\end{table}

On the fourth row the bulk matrix factorizes. With $\Gamma_5(u)=\Id+puE_{12}$,
\begin{equation}
R_5\big|_{r=-p}(u)=(\Gamma_5(u)\otimes\Gamma_5(-u))\,P,
\label{eq:C5-factorization}
\end{equation}
which says that on this line the Class-5 bulk matrix is nothing but the exchange
operator dressed by an entire twist on each leg. Substituting
\cref{eq:C5-factorization} into \cref{eq:re} and setting
\begin{equation}
K(u)=\Gamma_5(u)\Lambda(u)\Gamma_5(u)                    \label{eq:C5-twist}
\end{equation}
turns the reflection equation into $[\Lambda(u),\Lambda(v)]=0$, so
\cref{lem:centralizer} applies: modulo the scalar gauge $\Lambda=\Id+\psi B_0$
for one fixed traceless direction. Before the quotient of \cref{sec:quotient} the
direction ranges over the whole projective line of traceless matrices; after it,
four representatives suffice,
\begin{equation}
E_{12}+E_{21},\qquad E_{21},\qquad \sigma_z,\qquad E_{12},
\label{eq:C5-shapes}
\end{equation}
namely the two-sided, lower, semisimple and upper directions, together with the
scalar branch $\Gamma_5^2$. The polynomial family meets these at
\begin{equation}
\psi(u)=u,\qquad
B_0=\begin{pmatrix}b&c-2p\\d&-b\end{pmatrix},       \label{eq:C5-overlap}
\end{equation}
and in general the slope of \cref{eq:C5-twist} is $B=2pE_{12}+\psi'(0)B_0$: the
twist itself contributes a fixed upper-triangular piece to the first derivative.
At the origin of the parameter space $\Gamma_5=\Id$ and the three types are
scalar, semisimple $\Id+\psi\sigma_z$, and nilpotent $\Id+\psi E_{12}$.

\subsection{Class 6}
\label{subsec:c6}

Write $a=a_1$ and $Q=a_2$. Off the point where the bulk matrix becomes the
exchange operator, one quintic polynomial exhausts the boundaries:
\begin{equation}
K_6(u)=\begin{pmatrix}
1+bu-2aQdu^3&
cu-(4aQb+Q^2d)u^3+12a^2Q^2du^5\\
du&1-bu+2aQdu^3
\end{pmatrix}.                                      \label{eq:C6-poly}
\end{equation}
The odd powers are forced: the quadratic and quartic coefficients vanish
identically, so the boundary is an odd deformation of the identity in the two
diagonal entries and of $cu$ in the upper one. The partition is
\cref{tab:C6-cells}.

\begin{table}[t]
\centering\small
\caption{The complete Class-6 partition.}
\label{tab:C6-cells}
\begin{tabularx}{\textwidth}{>{\RaggedRight\arraybackslash}X >{\RaggedRight\arraybackslash}X}
\toprule
Condition & Complete family \\
\midrule
$a\ne0$, $Q\ne0$ & $K_6$, every slope \\
$a\ne0$, $Q=0$ & $K_6=\Id+uB$, every slope \\
$a=0$, $Q\ne0$ & $K_6$, every slope \\
$a=Q=0$ & every self-commuting family \\
\bottomrule
\end{tabularx}
\end{table}

On the last row the three types are again $\Id$, $\Id+\psi\sigma_z$ and
$\Id+\psi E_{12}$, with slope $B=\psi'(0)B_0$ for the chosen direction and zero
on the scalar branch. The specialization of \cref{eq:C6-poly} at the origin is
only the subfamily $\psi(u)=u$ and is not counted separately.

At $a=0$ with $Q\ne0$ the elimination can be run by hand and is worth displaying,
because it shows how the polynomial form is forced rather than assumed. In the
trace gauge, writing $p_0(u)=(k_1(u)-k_4(u))/2$, the system
\cref{eq:derivative-system} gives successively
\begin{equation}
k_3(u)=du,\qquad
k_2(u)=-Q^2du^3+2Qbu^2-2Qp_0(u)u+cu,\qquad
8Q\bigl(p_0(u)-bu\bigr)=0.
\label{eq:C6-elimination}
\end{equation}
The first equation fixes the lower entry, the second expresses the upper entry
through the remaining unknown $p_0$, and the third --- available because $Q\ne0$
--- forces $p_0(u)=bu$, which closes the system and reproduces
\cref{eq:C6-poly} at $a=0$.

\begin{proof}[Proof of \cref{thm:raw}]
\Cref{eq:derivative-system} is a necessary condition satisfied by every
admissible boundary, derived in \cref{sec:exhaustiveness} from four unrestricted
meromorphic entries. By \cref{lem:stratification} the exact radical, saturation
and rate-collision computations partition each
$\mathcal A_j\times\C^3$ into precisely the cells listed above, with no residual
component. On every rank-three cell the displayed normalized matrix spans the
kernel of $\mathcal C_{R,B}$, so it is the only candidate in the scalar gauge; on
every rank-two or rank-zero cell the specialized equations reduce, by the twists
displayed there, to a commutator identity whose solution set
\cref{lem:centralizer} exhausts. Sufficiency is separate from necessity, and is
supplied by substituting each displayed family back into the two-variable
identity \cref{eq:re}, with independent symbols for the values of every arbitrary
function at $u$ and at $v$; every residual vanishes identically. Finally
\cref{subsec:continuation} promotes the equality of normalized germs at the
origin to an identity of single-valued meromorphic functions on $\C$. The tables
are therefore valid and exhaustive on all six full parameter varieties.
\end{proof}

\section{The fixed-bulk quotient}
\label{sec:quotient}

\begin{theorem}[Orbits at fixed bulk data]
\label{thm:quotient}
For each family of \cref{thm:raw}, quotienting by \cref{eq:equivalence} gives
exactly the groups, orbit slices and stabilizers of this section. Every
transformation used preserves the same literal bulk matrix projectively in the
sense of \cref{eq:bulk-aut}; no parameter-changing similarity and no similarity
that is singular on a degeneration appears. The stabilizer jumps, the overlaps
between finite and arbitrary-function families, and the constant-bulk
degenerations are all included.
\end{theorem}

For Classes~1, 2, 4, 5 and 6, comparing determinants on both sides of
\cref{eq:bulk-aut} forces the bulk scalar $\varphi$ to be identically one: the
rescaling may not change the normalization of the bulk matrix at all. For
Class~3 the generator equation determines the scalar uniquely instead, and the
table below is for the literal source matrix.

Three families of similarity recur, and it is convenient to name them once:
\begin{equation}
U_\tau=\begin{pmatrix}1&\tau\\0&1\end{pmatrix},\qquad
S_{s,\tau}=\begin{pmatrix}1&\tau\\0&s\end{pmatrix},\qquad
D_s=\diag(1,s).
\label{eq:quotient-groups}
\end{equation}
They are the unipotent, Borel and torus subgroups of the upper-triangular
matrices; which of them survives at a given bulk point is what the tables record.
Their action on the slope is universal:
\begin{equation}
\begin{array}{c|c}
\text{additive}&(b,c,d)\mapsto(b+\tau d,\ c-2\tau b-\tau^2d,\ d)\\[1mm]
\text{affine}&(b,c,d)\mapsto(s(b+\tau d),\ c-2\tau b-\tau^2d,\ s^2d)\\[1mm]
\text{torus}&(b,c,d)\mapsto(sb,\ c,\ s^2d)\\[1mm]
\text{weighted Borel}&(b,c,d)\mapsto
(s^2(b+\tau d),\ s(c-2\tau b-\tau^2d),\ s^3d)
\end{array}                                             \label{eq:universal-actions}
\end{equation}
Each row of \cref{eq:universal-actions} is conjugation of the traceless first
derivative by the corresponding group, combined with the spectral rescaling where
one is allowed; the weights on $b$, $c$ and $d$ record how each entry scales.
Writing $\Delta_B=b^2+cd$ for the discriminant of the slope, which is minus its
determinant and therefore a conjugation invariant, the orbit slices are:
\begin{itemize}
\item additive: $(0,\Delta_B/d,d)$ for $d\ne0$, $(b,0,0)$ for $d=0$, $b\ne0$,
and $(0,c,0)$ for $b=d=0$;
\item affine: $(0,\Delta_B/d,1)$ for $d\ne0$, $(1,0,0)$ for $d=0$, $b\ne0$, and
$(0,c,0)$ for $b=d=0$;
\item torus: $(1,c,d/b^2)$ for $b\ne0$, $(0,c,1)$ for $b=0$, $d\ne0$, and
$(0,c,0)$ for $b=d=0$;
\item weighted Borel: $(0,\varsigma,1)$ for $d\ne0$ with
$\varsigma^3=\Delta_B^3/d^4$ modulo cube roots of unity; $(1,0,0)$ for $d=0$,
$b\ne0$; $(0,1,0)$ for $b=d=0$, $c\ne0$; and the origin.
\end{itemize}
For an arbitrary-function family $\Lambda=\Id+\psi B_0$ the group acts on the
fixed direction $B_0$ and on $\psi$ together, and the quotient is a list of
projective direction types each carrying a residual functional identity, as in
\cref{tab:borel-quotient,tab:normalizer-quotient}.

\begin{table}[t]
\centering\small
\caption{Quotient of an arbitrary-function fiber by the affine or weighted-Borel
group. The residual relation is the identification that survives on the free
function.}
\label{tab:borel-quotient}
\begin{tabularx}{\textwidth}{X c X}
\toprule
Direction cell & Representative & Residual relation \\
\midrule
$d\ne0$, $\Delta_B\ne0$ & $E_{12}+E_{21}$ & $\psi(u)\sim-\psi(-u)$ \\
$d\ne0$, $\Delta_B=0$ & $E_{21}$ & $\psi(u)\sim s\,\psi(su)$ \\
$d=0$, $b\ne0$ & $\sigma_z$ & $\psi(u)\sim \psi(su)$ \\
$b=d=0$, $c\ne0$ & $E_{12}$ & $\psi(u)\sim s^{-1}\psi(su)$ \\
\bottomrule
\end{tabularx}
\end{table}

\begin{table}[t]
\centering\small
\caption{Quotient of an arbitrary-function fiber by a torus normalizer, written
in the eigenbasis where
$\widetilde B_0=n_0\sigma_z+e_+E_{12}+e_-E_{21}$.}
\label{tab:normalizer-quotient}
\begin{tabularx}{\textwidth}{X X X}
\toprule
Direction cell & Representative & Residual relation \\
\midrule
$n_0\ne0$, $e_+e_-\ne0$ & $\sigma_z+E_{12}+\iota E_{21}$,
$\iota=e_+e_-/n_0^2$ & $\psi(u)\sim-\psi(-u)$ \\
$n_0\ne0$, exactly one of $e_\pm$ nonzero & $\sigma_z+E_{12}$ & none \\
$n_0\ne0$, $e_+=e_-=0$ & $\sigma_z$ & $\psi(u)\sim-\psi(-u)$ \\
$n_0=0$, $e_+e_-\ne0$ & $E_{12}+E_{21}$ & $\psi(u)\sim\pm\psi(\pm u)$ \\
$n_0=0$, exactly one of $e_\pm$ nonzero & $E_{12}$ & $\psi(u)\sim\varsigma\,\psi(u)$ \\
\bottomrule
\end{tabularx}
\end{table}

At a bulk point where the matrix is the exchange operator, the only nonscalar
directions are semisimple and nilpotent, with residual relations
$\psi(u)\sim\pm\psi(\lambda u)$ and $\psi(u)\sim\varsigma\,\psi(\lambda u)$
respectively, the two constants in the second being independent. In every row of
\cref{tab:borel-quotient,tab:normalizer-quotient} the stabilizer of a family is
exactly the subgroup for which the displayed relation holds, and the scalar
branch is fixed by the whole group.

\subsection{Classes 1 and 2}

For Class~1 set $Z_1=\{a_4=a_1=t,\ a_2=a_3=-t\}$ and
$\ell_1=a_3+a_4-a_1-a_2$; on $Z_1$ put
$N_1=\begin{psmallmatrix}1&t/a_5\\0&1\end{psmallmatrix}$ and write
$W=\begin{psmallmatrix}0&1\\1&0\end{psmallmatrix}$ for the swap. The complete
group table is \cref{tab:C1-aut}.

\begin{table}[t]
\centering\footnotesize
\caption{Class-1 fixed-bulk automorphisms, as pairs $(\lambda,[S])$.}
\label{tab:C1-aut}
\begin{tabularx}{\textwidth}{>{\RaggedRight\arraybackslash}X >{\RaggedRight\arraybackslash}X}
\toprule
Bulk condition & Complete group \\
\midrule
$a_5\ne0$, on $Z_1$ &
$(1,[N_1D_sN_1^{-1}])$ and $(-1,[N_1WD_sN_1^{-1}])$ \\
$a_5\ne0$, off $Z_1$ & trivial \\
$a_5=0$, all remaining parameters zero & every $\lambda$ and all of $PGL_2$ \\
$a_5=0$, some parameter nonzero, $\ell_1=0$ & $(s,[S_{s,\tau}])$ \\
$a_5=0$, $\ell_1\ne0$ & $(s,[D_s])$ \\
\bottomrule
\end{tabularx}
\end{table}

For Class~2, on $a_3=\varepsilon=0$ with $a_1\ne0$ put
\begin{equation}
N_2=\begin{pmatrix}1&a_2/a_1\\0&1\end{pmatrix},\qquad
W_2=\begin{pmatrix}1&-2a_2/a_1\\0&-1\end{pmatrix},
\label{eq:C2-conjugators}
\end{equation}
the conjugator that straightens the bulk matrix and the involution it produces.
The group table is \cref{tab:C2-aut}.

\begin{table}[t]
\centering\footnotesize
\caption{Class-2 fixed-bulk automorphisms, as pairs $(\lambda,[S])$.}
\label{tab:C2-aut}
\begin{tabularx}{\textwidth}{>{\RaggedRight\arraybackslash}X >{\RaggedRight\arraybackslash}X}
\toprule
Bulk condition & Complete group \\
\midrule
$a_1\ne0$, $a_3=\varepsilon=a_5=0$ & two $N_2$-conjugated normalizer components, with $\lambda=\pm1$ \\
$a_1\ne0$, $a_3=\varepsilon=0$, $a_5\ne0$ & $(1,[\Id])$ and $(1,[W_2])$ \\
$a_1\ne0$, all other points & trivial \\
$a_1=0$, all remaining parameters zero & every $\lambda$ and all of $PGL_2$ \\
$a_1=0$, $a_2=a_3=a_4=0$, $a_5\ne0$ & $(s^2,[S_{s,\tau}])$ \\
$a_1=0$, nonzero linear part, $a_5=0$ & $(s,[S_{s,\tau}])$ \\
$a_1=0$, nonzero linear part, $a_5\ne0$ & $(1,[U_\tau])$ \\
\bottomrule
\end{tabularx}
\end{table}

The routing is as follows. On the first Class-1 row use
\cref{tab:normalizer-quotient} with $\widetilde B_0=N_1^{-1}B_0N_1$; the second
row has trivial group, so every raw family there is already a normal form; the
third uses the constant-bulk rules; the fourth uses
\cref{tab:borel-quotient} on $\mathcal E_1$ and the affine slices off it; the
fifth uses the torus slices, with two exceptional rows,
\begin{align}
K&=\Id+\psi(u)E_{12},&
\psi(u)&\sim s^{-1}\psi(su),
&&a_1=a_4=0,\ (a_2,a_3)\ne(0,0),\ a_2\ne a_3,\label{eq:C1-upper-orbit}\\
K&=\diag(1,\psi(u)),&
\psi(u)&\sim \psi(su),
&&a_2=a_3=0,\ (a_1,a_4)\ne(0,0),\ a_1\ne a_4,\label{eq:C1-diag-orbit}
\end{align}
whose stabilizers are exactly the sets of $s$ satisfying the displayed functional
identity, the constant members $\psi=0$ and $\psi=1$ being fixed by the whole
torus.

For Class~2, on the first row use \cref{tab:normalizer-quotient} with
$\widetilde B_0=N_2^{-1}B_0N_2$. On the second row put
$\widehat c=c+2a_2b/a_1$; the involution fixes $b$ and sends
$\widehat c\mapsto-\widehat c$, so the unique rows are labelled by
$(b,\widehat c^{\,2})$, with stabilizer of order two exactly at $\widehat c=0$.
At $b=-a_1$, if $\psi$ is the arbitrary upper entry, set
\begin{equation}
\widetilde\psi(u)=\psi(u)-\frac{a_2}{a_1}\left(e^{a_1u}-e^{-a_1u}\right);
\label{eq:C2-shifted-function}
\end{equation}
subtracting this fixed particular solution makes the residual relation homogeneous,
$\widetilde\psi\sim-\widetilde\psi$, whose stabilizer has order two exactly when
$\widetilde\psi=0$. The remaining rows use the constant-bulk rules, the
weighted-Borel slices, the Borel table on $\varepsilon=0$, the affine slices on
$a_3=0$, $\varepsilon\ne0$, and the additive slices otherwise; at $b=d=0$ with
$a_2=a_3=a_4=0$, $a_5\ne0$ the orbit is
\begin{equation}
K=\Id+\psi(u)E_{12},\qquad \psi(u)\sim s^{-1}\psi(s^2u),       \label{eq:C2-weighted-upper}
\end{equation}
and every $\tau$ acts trivially there.

\subsection{Class 3}

With $(m,n,g)$ as in \cref{eq:C3-vars}, the fixed-bulk groups are those of
\cref{tab:C3-aut}: a nonzero coupling $g$ rigidifies the bulk matrix completely,
and only at $g=0$ does a torus survive.

\begin{table}[t]
\centering\small
\caption{Class-3 fixed-bulk automorphisms.}
\label{tab:C3-aut}
\begin{tabularx}{\textwidth}{X X}
\toprule
Bulk stratum & Complete set of pairs $(\lambda,[S])$ \\
\midrule
$m=n=0$, any $g$ & every $\lambda\ne0$ and every $[S]\in PGL_2$ \\
$g\ne0$, $(m,n)\ne(0,0)$ & trivial \\
$g=0$, $(m,n)\ne(0,0)$, $m\ne\pm n$ & $\lambda=1$, $S=D_s$ \\
$g=0$, $m=n\ne0$ & $\lambda=1$, $S$ diagonal or antidiagonal \\
$g=0$, $m=-n\ne0$ & $\lambda=1$ with $S$ diagonal, or $\lambda=-1$ with $S$ antidiagonal \\
\bottomrule
\end{tabularx}
\end{table}

On the diagonal-torus cell the slope transforms as $(b,c,d)\mapsto(b,c/s,sd)$, so
the two-sided directions have slice $(b,1,cd)$ while the upper and lower one-sided
directions have $(b,1,0)$ and $(b,0,1)$. At $m=n$ the extra antidiagonal element
sends $(b,c,d)\mapsto(-b,d/s,sc)$, identifying $b$ with $-b$ on the two-sided
slice and sending a diagonal family to its reciprocal.

On $g=0$, $m=-n\ne0$, the homogeneous family \cref{eq:C3-opposite} has two
projective direction cells. For $\gamma_1\gamma_2\ne0$ take $\gamma_1=\gamma_2=1$;
the residual equivalence is generated by
\begin{equation}
(\psi,\nu)\mapsto(-\psi,-\nu),
\qquad
(\psi(u),\nu)\mapsto
\left(\frac{\theta(u)^2\psi(-u)}{1+\nu\,\theta(u)\,\psi(-u)},\ -\nu\right),
\label{eq:C3-residual}
\end{equation}
the second of which is the reflection $u\mapsto-u$ acting through a Möbius
transformation on the free function. For $\gamma_1\gamma_2=0$ the upper and lower
charts merge; taking $\gamma_1=1$, $\gamma_2=0$ and using
$(\psi,\nu)\mapsto(\psi/s,s\nu)$ leaves the two slices $\nu=0$ and $\nu=1$. A
diagonal family obeys $\psi(u)\sim\psi(-u)^{-1}$.

\subsection{Class 4}

The eight fixed-bulk strata are those of \cref{tab:C4-aut}.

\begin{table}[t]
\centering\footnotesize
\caption{Class-4 fixed-bulk automorphisms.}
\label{tab:C4-aut}
\begin{tabularx}{\textwidth}{>{\RaggedRight\arraybackslash}X >{\RaggedRight\arraybackslash}X}
\toprule
Bulk condition & Complete group \\
\midrule
$a_1\ne0$, $a_2+a_4\ne0$ & trivial \\
$a_1\ne0$, $a_2+a_4=0$, $a_1a_3+a_2a_4\ne0$ & the identity and $\begin{psmallmatrix}1&a_2/a_1\\0&-1\end{psmallmatrix}$, with $\lambda=1$ \\
$a_1\ne0$, $a_2+a_4=0$, $a_1a_3+a_2a_4=0$ & a conjugated diagonal-torus normalizer, $\lambda=1$ \\
$a_1=0$, $a_2\ne a_4$ & a conjugated torus, spectral scale $\lambda=s$ \\
$a_1=0$, $a_2=a_4\ne0$, $a_3\ne0$ & the additive group, $\lambda=1$ \\
$a_1=0$, $a_2=a_4\ne0$, $a_3=0$ & the affine group, spectral scale $\lambda=s$ \\
$a_1=0$, $a_2=a_4=0$, $a_3\ne0$ & the weighted Borel group, spectral scale $\lambda=s^2$ \\
zero bulk & every $\lambda$ and all of $PGL_2$ \\
\bottomrule
\end{tabularx}
\end{table}

On the second row, the shifted slope coordinates
\begin{equation}
b_1=b_0+\tfrac{A}{2}d_0,\qquad
c_1=c_0-Ab_0-\tfrac{A^2}{4}d_0,\qquad d_1=d_0
\label{eq:C4-shift}
\end{equation}
diagonalize the involution, which acts as $(b_1,c_1,d_1)\mapsto(b_1,-c_1,-d_1)$;
the shift is exactly the conjugation that removes the off-diagonal part of the
automorphism. On the third row the torus acts by
$(b_1,c_1,d_1)\mapsto(b_1,c_1/s,sd_1)$ and its Weyl coset exchanges $c_1$ and
$d_1$ while sending $b_1\mapsto-b_1$. On the fourth, after conjugation by
$\Id+\frac{a_3}{2(a_2-a_4)}E_{12}$, the action is
$(b_1,c_1,d_1)\mapsto(sb_1,c_1,s^2d_1)$. The remaining nonzero-bulk rows use the
additive, affine and weighted-Borel actions of \cref{eq:universal-actions}, and
the zero-bulk row uses the three constant-bulk types.

\subsection{Classes 5 and 6}

The Class-5 group changes on four cells:
\begin{equation}
\begin{array}{c|c}
a\ne0,\ (p,r)\ne(0,0)&(1,[U_\tau])\\
a\ne0,\ p=r=0&(1,PGL_2)\\
a=0,\ (p,r)\ne(0,0)&(s,[S_{s,\tau}])\\
a=p=r=0&(\lambda,PGL_2).
\end{array}                                           \label{eq:C5-aut}
\end{equation}
The pattern is that switching off a bulk parameter enlarges the group, from a
single unipotent direction at a generic point to the whole of $PGL_2$ at the
origin. On the polynomial cells the action is the additive or affine row of
\cref{eq:universal-actions}; on the cell where the bulk matrix is rational and
the group is all of $PGL_2$, the invariant is $\Delta_B=b^2+cd$ and the orbits
are the scalar, nonzero nilpotent and semisimple types. On the twisted cell the
four directions \cref{eq:C5-shapes} and their residual relations are those of
\cref{tab:borel-quotient}, merging at the origin into the semisimple and
nilpotent constant-bulk types.

For Class~6 the four groups are
\begin{equation}
\begin{array}{c|c}
a\ne0,\ Q\ne0&(1,[\Id])\\
a\ne0,\ Q=0&(1,PGL_2)\\
a=0,\ Q\ne0&(s^{-1},[D_s])\\
a=Q=0&\text{every }\lambda\text{ and all of }PGL_2 .
\end{array}                                           \label{eq:C6-aut}
\end{equation}
On the torus cell the slope transforms as $(b,c,d)\mapsto(b/s,c,d/s^2)$, with
slices $(1,c,d/b^2)$ for $b\ne0$, $(0,c,1)$ for $b=0$, $d\ne0$, and $(0,c,0)$ for
$b=d=0$. The rational and constant fibers carry the same $PGL_2$ orbit types by
Jordan type as above. These rules also state every degeneration: a group may
enlarge as a bulk parameter vanishes, but the enlargement is never used to
identify two different bulk points, which would be a change of model rather than
a change of representative.

\begin{proof}[Proof of \cref{thm:quotient}]
Differentiating \cref{eq:bulk-aut} at the origin gives an exact equation for the
generator of the similarity in the tensor square, linear in the unknown entries of
$S$ and in $\lambda$. Comparing determinants and spectra on the two sides
restricts the spectral scale, and solving the generator equation on both
projective charts of $S$, with the relevant denominators saturated, yields exactly
the strata tabulated above; the exponential, hyperbolic or polynomial form of each
bulk matrix makes the infinitesimal condition sufficient, and the one remaining
Class-4 covariant is verified directly. Conjugating the traceless first derivative
by the resulting groups gives \cref{eq:universal-actions} and the shifted Class-3
and Class-4 variants; solving the stabilizer equations gives the displayed slices.
On an arbitrary-function fiber, undoing the twist reduces the action to that of
the normalizer, Borel or $PGL_2$ on a single fixed traceless direction, whose
orbit types are the semisimple, nilpotent, two-sided and one-sided shapes listed.
Direct substitution confirms every inclusion and every overlap between a finite
and an arbitrary-function family. No further identification at fixed bulk data
remains.
\end{proof}

\section{Crossing, left boundaries, and open chains}
\label{sec:crossing}

\subsection{The crossing datum and what it is for}

An open chain needs a boundary at each end. The standard way to obtain the second
from the first is a crossing relation: constants $M\in GL_2(\C)$ and
$\rho\in\C$ and a nonzero meromorphic $\chi$ with
\begin{equation}
R_{12}^{t_1}(u)\,M_1\,R_{21}^{t_1}(-u-2\rho)\,M_1^{-1}
=\chi(u)\,I_4.                                         \label{eq:crossing}
\end{equation}
\Cref{eq:crossing} is the algebraic form of particle--antiparticle exchange: it
says the first-leg partial transpose of the bulk matrix is inverted, up to a
scalar and a fixed conjugation, by the same matrix at a reflected and shifted
rapidity. Given such a datum, a right boundary $\widetilde K_R$ prescribes the
left boundary
\begin{equation}
K_L(u)=\widetilde K_R(-u-\rho)^{t}M.                  \label{eq:left-map}
\end{equation}
For a chain of $L$ sites, write the two monodromies and the double-row transfer
matrix as
\begin{equation}
T_a(u)=R_{aL}(u)\cdots R_{a1}(u),\qquad
\widehat T_a(u)=R_{1a}(u)\cdots R_{La}(u),
\label{eq:monodromies}
\end{equation}
\begin{equation}
t_L(u)=\Tr_a\left[K_{L,a}(u)\,T_a(u)\,K_{R,a}(u)\,
\widehat T_a(u)\right].                             \label{eq:transfer}
\end{equation}
The auxiliary line crosses all $L$ sites, reflects at the right end through
$K_R$, crosses back, and is closed at the left end by $K_L$ under the trace over
the auxiliary space; the commuting tower of the open chain is generated by
$t_L$. We call the left map \emph{usable} on a cell when
$[t_L(u),t_L(v)]=0$ for every $L\ge2$ and every ordered pair
$(\widetilde K_R,K_R)$ of independently chosen classified right boundaries --- not
merely for one boundary, and not merely for short chains.

\subsection{Classes 1--4 admit no crossing datum}

\begin{proposition}
\label{prop:no-crossing}
For every $j\in\{1,2,3,4\}$ and every $\mathbf a\in\mathcal A_j$ there is no
triple $(M,\rho,\chi)$ with $M$ invertible and $\chi$ not identically zero
satisfying \cref{eq:crossing}.
\end{proposition}

\begin{proof}
Up to an overall scalar, each of \cref{eq:R1,eq:R2,eq:R3,eq:R4} has the pattern
\begin{equation}
R(u)=\begin{pmatrix}
1&\varrho_{12}(u)&\varrho_{13}(u)&\varrho_{14}(u)\\
0&0&\varrho_{23}(u)&\varrho_{24}(u)\\
0&\varrho_{32}(u)&0&\varrho_{34}(u)\\
0&0&0&1
\end{pmatrix},
\label{eq:pattern}
\end{equation}
that is, a matrix whose first column and last row are those of the identity and
whose middle block is antidiagonal. Its first-leg partial transpose is
\begin{equation}
R^{t_1}(u)=\begin{pmatrix}
1&\varrho_{12}(u)&0&\varrho_{32}(u)\\
0&0&0&0\\
\varrho_{13}(u)&\varrho_{14}(u)&0&\varrho_{34}(u)\\
\varrho_{23}(u)&\varrho_{24}(u)&0&1
\end{pmatrix}.                                      \label{eq:t1-rank}
\end{equation}
The second row and the third column of \cref{eq:t1-rank} vanish identically, so
$R^{t_1}(u)$ has rank at most three at every rapidity and every bulk point.
Suppose \cref{eq:crossing} held. Choose $u$ away from the poles and zeros of the
finitely many meromorphic functions involved and of $\chi$. The left side is a
product of $R^{t_1}(u)$ with invertible matrices, hence of rank at most three,
while $\chi(u)I_4$ has rank four. The contradiction is uniform in the bulk
parameters, so no crossing datum exists anywhere on those four varieties.
\end{proof}

No left map, transfer matrix, or boundary Hamiltonian obligation of the form
\cref{eq:left-map} is therefore triggered for Classes~1--4. This is a genuine
structural statement about these models and not a failure of search: the
partial transpose of \cref{eq:pattern} is singular by inspection.

\subsection{Classes 5 and 6: crossing exists exactly at nonzero rate}

\begin{proposition}
\label{prop:crossing-56}
For Class~5 with $(a,p,r)=(a_1,a_2,a_3)$, a crossing datum exists if and only if
$a\ne0$, and is then unique up to a scalar multiple of $M$:
\begin{equation}
M=\begin{pmatrix}1&0\\(p-r)/(2a)&1\end{pmatrix},\qquad
\rho=\frac1{2a}.                                    \label{eq:C5-crossing}
\end{equation}
For Class~6 with $(a,Q)=(a_1,a_2)$, a crossing datum exists if and only if
$a\ne0$, uniquely with
\begin{equation}
M=\Id,\qquad \rho=\frac1{2a}.                       \label{eq:C6-crossing}
\end{equation}
For both source-normalized matrices the crossing scalar is
\begin{equation}
\chi(u)=4au(au-1)(au+1)(au+2).                     \label{eq:chi}
\end{equation}
\end{proposition}

The shift $\rho=1/(2a)$ is the pole of the bulk prefactor: crossing reflects the
rapidity about the point where the bulk matrix degenerates. \Cref{eq:chi} shows
that the crossing scalar vanishes at the origin and at three further points, all
of them fixed by the rate. Existence and uniqueness follow from exact elimination
of the coefficients of \cref{eq:crossing} under the inequation $\det M\ne0$; at
$a=0$ the elimination has no solution, and on the remaining constant and twisted
fibers the rank of $R^{t_1}$ supplies the same obstruction as in
\cref{prop:no-crossing}.

Crossing existence is not enough for usability.

\begin{proposition}
\label{prop:not-usable}
On Class~6 with $aQ\ne0$ the left map \cref{eq:left-map} is not usable. Taking
$K_R=\Id$ and the independent seed with slope $(b,c,d)=(1,0,0)$, and expanding
$t_2(u)=\sum_{k\ge0}t^{(k)}u^k$, one has
\begin{equation}
\bigl[t^{(1)},t^{(2)}\bigr]_{1,2}=64Q^2\ne0.        \label{eq:C6-no-go}
\end{equation}
On Class~5 with $a\ne0$ and $p\ne r$ the ideal generated by the divided
commutator coefficients, saturated by the inequation $a\ne0$, is the unit ideal;
and on $p=r\ne0$ the independent upper seed gives
\begin{equation}
\bigl[t^{(1)},t^{(2)}\bigr]_{0,2}=16r^2\ne0.        \label{eq:C5-no-go}
\end{equation}
\end{proposition}

Each display is a single matrix entry of a commutator that would have to vanish,
computed for a chain of two sites and exhibited as a nonzero number. Together
with \cref{prop:crossing-56} they leave exactly the rational cells
\begin{equation}
a\ne0,\quad p=r=0\ \text{ (Class 5)},\qquad
a\ne0,\quad Q=0\ \text{ (Class 6)}.              \label{eq:rational-cell}
\end{equation}

\subsection{All-length commutativity on the rational cells}

On \cref{eq:rational-cell} both classes reduce to the same rational bulk matrix
\begin{equation}
R(u)=(1-au)(P+2au\,I_4),                              \label{eq:rational-R}
\end{equation}
the exchange operator plus a multiple of the identity, times a scalar; this is
the isotropic point of both families. There every classified right boundary is
$\Id+uB_R$ with $\tr B_R=0$, and for an independently chosen seed with slope
$\widetilde B$ the prescribed left boundary is
\begin{equation}
K_L^{\mathrm{can}}(u)=\Id-(u+\rho)\widetilde B^{\,t}.        \label{eq:rational-left}
\end{equation}
The left boundary is thus the seed reflected about $-\rho$ and transposed, which
is \cref{eq:left-map} written out for a linear boundary.

\begin{theorem}
\label{thm:rational-cell}
On the cells \cref{eq:rational-cell}, for every $L\ge2$ and every ordered pair of
classified boundaries, $[t_L(u),t_L(v)]=0$.
\end{theorem}

\begin{proof}
For \cref{eq:rational-R} the Yang--Baxter equation, regularity, the permutation
and full-transpose symmetries, scalar unitarity, and the crossing identity
\cref{eq:crossing} all hold, and the dual reflection equation holds for arbitrary
and independent $B_R$ and $\widetilde B$ by direct substitution. The
reflection-algebra argument of \cite{sklyanin1988}, which builds the commutator
of two double rows by moving one train of bulk matrices through the other and
uses the two reflection equations to reverse the order at the ends, then applies
verbatim and gives the claim for every $L$.
\end{proof}

Literal representatives, rather than canonically normalized ones, must be
retained, because the scalar gauge is part of the classification. Writing
$\widetilde K_R(u)=\widetilde f(u)(\Id+u\widetilde B)$ and
$K_R(u)=f_R(u)(\Id+uB_R)$, linearity of \cref{eq:transfer} in each boundary gives
\begin{equation}
t_L(u)=\widetilde f(-u-\rho)\,f_R(u)\,t_L^{\mathrm{can}}(u).
\label{eq:literal-transfer}
\end{equation}
A change of gauge therefore multiplies the whole transfer matrix by a scalar
function and cannot affect commutativity; but it can move the expansion point.
Exactly three cases occur at $u=0$: if $\widetilde f(-\rho)$ is finite and
nonzero, $t_L(0)$ is invertible; if it vanishes, $t_L(0)=0$; and if it has a
pole, $K_L$ and $t_L$ have a pole at the origin of the same order. The admissible
gauges $1+2au$ and $(1+2au)^{-1}$ realize the last two cases, and no logarithmic
derivative is taken there.

\subsection{The open-chain Hamiltonian and an exact normalization obstruction}

For canonically normalized pairs on the rational cells,
$t_L^{\mathrm{can}}(0)=2\,I_{V^{\otimes L}}$ and the two-site density is
\begin{equation}
h=PR'(0)=2aP-a\,I_4
=a\,(\sigma_x\otimes\sigma_x+\sigma_y\otimes\sigma_y+\sigma_z\otimes\sigma_z),
\label{eq:rational-density}
\end{equation}
the isotropic Heisenberg interaction with coupling $a$. Differentiating the
permutation train at the origin gives
\begin{equation}
\mathcal H_L=t_L'(0)\,t_L(0)^{-1}
=2\sum_{n=1}^{L-1}h_{n,n+1}+(B_R)_1-(\widetilde B^{\,t})_L,
             \label{eq:H-exact}
\end{equation}
modulo a multiple of the identity: the open chain is the bulk sum with one
one-site term at each end, the right end carrying the slope of the right boundary
and the left end the transposed slope of the seed. Finite nonzero literal gauges
change only the additive scalar.

The bulk sum in \cref{eq:H-exact} carries a coefficient two rather than one, and
this cannot be repaired by adjusting the endpoints.

\begin{proposition}
\label{prop:normalization}
There is no choice of one-site operators at sites $1$ and $L$, and no multiple of
the identity, for which \cref{eq:H-exact} equals $\sum_n h_{n,n+1}$ plus those
one-site terms.
\end{proposition}

\begin{proof}
Take $L=4$ and expand both sides in the basis of tensor products of Pauli
matrices. The coefficient of
$\Id\otimes\sigma_x\otimes\sigma_x\otimes\Id$ in the excess $h_{23}$ is $a\ne0$,
while an operator supported at site $1$ or site $4$ and a multiple of the identity
both have coefficient zero there.
\end{proof}

The factor two is the record of the two passages of the auxiliary line through
every bulk bond in \cref{eq:transfer}. It is a property of the unhalved
logarithmic derivative used here, not a missing boundary term; halving it, as is
customary, restores the conventional coefficient.

\section{A constant left boundary for every Class-6 boundary}
\label{sec:class6}

\Cref{prop:not-usable} says that no single crossing-induced left boundary works
against every independently chosen right boundary off the rational cells. That is
a stronger demand than an open chain makes. To build one commuting tower it is
enough to exhibit, for each right boundary, one compatible left boundary. For
Class~6 there is one, and it is the same for all of them.

\begin{theorem}
\label{thm:class6}
Let $R_6$ be the source-normalized Class-6 matrix at any
$(a,Q)\in\C^2$, let $K_R$ be any admissible boundary classified in
\cref{subsec:c6}, including the arbitrary-function families at $a=Q=0$, and let
$K_L(u)=\Id$. Then for every $L\ge2$ the double-row transfer matrix
\begin{equation}
t_L(u)=\Tr_a\!\left[T_a(u)\,K_{R,a}(u)\,\widehat T_a(u)\right]
\label{eq:C6-transfer}
\end{equation}
satisfies $[t_L(u),t_L(v)]=0$ as an identity of meromorphic functions. If
$K_R(0)=\kappa\Id$ with $\kappa\ne0$, then $t_L(0)=2\kappa I$ is invertible and
\begin{equation}
t_L'(0)\,t_L(0)^{-1}
=2\sum_{n=1}^{L-1}h_{n,n+1}
 +\bigl(\kappa^{-1}K_R'(0)\bigr)_1,
\qquad h=P R_6'(0).                                \label{eq:C6-hamiltonian}
\end{equation}
\end{theorem}

\Cref{eq:C6-hamiltonian} says that the open Class-6 chain built this way is the
bulk chain with a single one-site term at the boundary site, given by the
derivative of the boundary matrix, and with no term at the auxiliary end at all.
Note what the theorem does not say: it exhibits one compatible left boundary for
each right boundary and does not assert that an arbitrary crossing-induced left
boundary works at $aQ\ne0$, which \cref{prop:not-usable} refutes.

\subsection{The Class-6 data}

Cancel the removable denominator in the middle block of \cref{eq:R6}, keeping
the source scalar; this gives
\begin{equation}
R_6(u)=(1-au)
\begin{pmatrix}
1+2au&Qu(1+2au)&Qu(1+2au)&-Q^2u^2(1+2au)^2\\
0&2au&1&-Qu(1+2au)\\
0&1&2au&-Qu(1+2au)\\
0&0&0&1+2au
\end{pmatrix},
\label{eq:C6-explicit}
\end{equation}
a matrix polynomial in the rapidity of degree four. Two features are used below
and both are visible in \cref{eq:C6-explicit}: the middle block is symmetric, so
\begin{equation}
R_{21}(u)=R_{12}(u),
\label{eq:C6-PT}
\end{equation}
meaning that Class~6 is one of the families for which the two orderings of
\cref{eq:re} agree; and braided unitarity holds in the sharp form
\begin{equation}
R_{12}(u)R_{21}(-u)=\varpi(u)\,I_4,
\qquad
\varpi(u)=(au-1)(au+1)(2au-1)(2au+1),
\label{eq:C6-unitarity}
\end{equation}
so exchanging two excitations and exchanging them back is the identity up to an
explicit quartic scalar. The determinants
\begin{equation}
\det R_6(u)=(au-1)^4(2au-1)(2au+1)^3,\qquad
\det R_6^{t_1}(u)=16a^3u^3(au-1)^4(au+1)
\label{eq:C6-determinants}
\end{equation}
show that $R_6$ and its first-leg partial transpose are invertible as meromorphic
matrix functions exactly when $a\ne0$; at $a=0$ the partial transpose is
identically singular, which is why that stratum needs a separate argument.

\subsection{The generic stratum}

\begin{proof}[Proof of \cref{thm:class6} for $a\ne0$]
Set $\widetilde R(u)=\bigl((R(u)^{t_1})^{-1}\bigr)^{t_1}$, which by
\cref{eq:C6-determinants} is a well-defined meromorphic matrix function for
$a\ne0$. The boundary transfer-matrix theorem of \cite[Thm.~2.4]{Vlaar_2015bt}
derives commutativity of two boundary transfer matrices under three hypotheses:
that $R$ and $R^{t_1}$ are invertible as meromorphic matrix functions, which
\cref{eq:C6-determinants} supplies for $a\ne0$; the right boundary-monodromy
reflection equation; and an operator-valued dual reflection equation for the
left boundary. For a one-dimensional left space and
$K_L=\Id$ the latter reduces to
\begin{equation}
R_{21}(u-v)^{-t}\,\widetilde R_{12}(u+v)^{t}
=\widetilde R_{21}(u+v)^{t}\,R_{12}(u-v)^{-t},
\label{eq:C6-dual}
\end{equation}
where $X^{-t}$ abbreviates $(X^{-1})^t$. Direct substitution of
\cref{eq:C6-explicit} into \cref{eq:C6-dual} gives a zero $4\times4$ residual in
the unrestricted symbols $a$, $Q$, $u-v$ and $u+v$: the constant left boundary
satisfies the dual equation identically, with no crossing relation, no $PT$
hypothesis, and no restriction on $Q$ entering.

It remains to connect the theorem's dressed monodromy to \cref{eq:C6-transfer}.
The permutation symmetry \cref{eq:C6-PT} and the local Yang--Baxter equation give
the global Yang--Baxter relation for $\widehat T_a$, and braided unitarity
\cref{eq:C6-unitarity} gives, meromorphically,
\begin{equation}
\widehat T_a(-u)^{-1}=\varpi(u)^{-L}\,T_a(u),
\label{eq:C6-unitarity-monodromy}
\end{equation}
that is, reversing the auxiliary line and inverting it returns the other
monodromy up to the $L$-th power of the unitarity scalar. Dressing the right
boundary in the standard way therefore produces
\begin{equation}
\mathcal K_a(u)
=\widehat T_a(-u)^{-1}K_{R,a}(u)\widehat T_a(u)
=\varpi(u)^{-L}\,T_a(u)K_{R,a}(u)\widehat T_a(u),
\label{eq:C6-dressing}
\end{equation}
which obeys the right boundary-monodromy reflection equation because $K_R$ obeys
\cref{eq:re}. Applying the theorem to \cref{eq:C6-dual,eq:C6-dressing} gives a
commuting family $\Tr_a\mathcal K_a(u)=\varpi(u)^{-L}t_L(u)$, and multiplying by
the scalar meromorphic function $\varpi(u)^{L}$ leaves the commutator unchanged.
Hence $[t_L(u),t_L(v)]=0$ for every $L$. Multiplying $K_R$ by an admissible
scalar gauge multiplies $t_L$ by the same function, so the conclusion holds for
every literal representative and not only for the trace-normalized formula.
\end{proof}

\begin{remark}
\label{rem:dual-display}
The scalar specialization of the dual reflection equation as displayed in
\cite[Eq.~(2.13)]{Vlaar_2015bt} ends in $R_{21}^{-t}$, whereas transposing the
operator equation of \cite[Thm.~2.4]{Vlaar_2015bt} in full ends in
$R_{21}^{-1}$, which is the form used in that paper's own appendix. Substituting
\cref{eq:C6-explicit} separates the two: the operator equation and the form
ending in $R_{21}^{-1}$ both give a zero residual, while the printed
specialization leaves fifteen nonzero entries for generic $Q$. The proof above
uses the operator equation directly and so assumes nothing about the printed
specialization.
\end{remark}

\subsection{The singular strata}

\begin{proof}[Proof of \cref{thm:class6} at $a=0$]
Fix $L$ and take first $Q\ne0$. For the classified boundary \cref{eq:C6-poly},
every entry of the commutator
\begin{equation}
C_L(a,Q,b,c,d;u,v)=[t_L(u),t_L(v)]
\label{eq:C6-commutator-poly}
\end{equation}
is a polynomial in all seven displayed variables, because
\cref{eq:C6-explicit,eq:C6-poly} are polynomial and \cref{eq:C6-transfer} is a
finite product and trace. The previous subsection proves that this polynomial
vanishes for every $a\ne0$ with $Q,b,c,d,u,v$ unrestricted. The complement of
$a=0$ is Zariski dense in $\C$, so every entry of $C_L$ is the zero polynomial,
and evaluating at $a=0$ proves commutativity on that stratum. Since $L$ was
arbitrary, this is an all-length continuation and not an inference from
finite-length checks.

At $a=Q=0$ the bulk matrix is $P$ and \cref{eq:re} reads $[K_R(u),K_R(v)]=0$, so
the admissible boundaries are the families of \cref{lem:centralizer}. There
$\widehat T_a(0)=T_a(0)^{-1}$, and conjugating by the permutation train moves the
auxiliary boundary onto the first site,
\begin{equation}
T_a(u)K_{R,a}(u)\widehat T_a(u)=\bigl(K_R(u)\bigr)_1,
\label{eq:C6-P-reduction}
\end{equation}
after which the auxiliary factor is idle and the trace contributes a factor two:
\begin{equation}
t_L(u)=2\bigl(K_R(u)\bigr)_1.
\label{eq:C6-P-transfer}
\end{equation}
The transfer matrix is thus the boundary matrix itself, acting on one site. Its
commutator is $4[K_R(u),K_R(v)]_1=0$, which covers the scalar, semisimple and
nilpotent arbitrary-function families alike, including those whose first nonzero
term occurs at arbitrarily high order.
\end{proof}

\subsection{The logarithmic derivative}

\begin{proof}[Proof of \cref{eq:C6-hamiltonian}]
For \cref{eq:C6-explicit} the two-site density is
\begin{equation}
h=P R_6'(0)=
\begin{pmatrix}
a&Q&Q&0\\
0&-a&2a&-Q\\
0&2a&-a&-Q\\
0&0&0&a
\end{pmatrix}
=a\,(\sigma_x\otimes\sigma_x+\sigma_y\otimes\sigma_y+\sigma_z\otimes\sigma_z)
+Q\,(\sigma_z\otimes E_{12}+E_{12}\otimes\sigma_z),
\label{eq:C6-density}
\end{equation}
an isotropic exchange of strength $a$ plus a non-Hermitian term of strength $Q$
that raises one spin while measuring the other. Both partial traces of $h$
vanish, which is the fact used at the end of this proof.

At $u=0$ the two permutation monodromies cancel and $K_R(0)=\kappa\Id$, so
$t_L(0)=2\kappa I$. Differentiating \cref{eq:C6-transfer} at the origin before
taking the auxiliary trace separates the terms into three groups. Differentiating
the boundary conjugates $K_R'(0)$ from the auxiliary factor onto site~$1$ and,
after the trace, contributes $2(K_R'(0))_1$. For each internal bond $(n,n+1)$,
differentiating the forward and the reflected copy of $R$ gives two equal
contributions, together $4\kappa h_{n,n+1}$. The two remaining terms, in which
the differentiated bulk matrix sits at an end of the train adjacent to the
auxiliary space, reduce after the trace to partial traces of $h$ and vanish.
Hence
\begin{equation}
t_L'(0)=4\kappa\sum_{n=1}^{L-1}h_{n,n+1}+2\bigl(K_R'(0)\bigr)_1,
\label{eq:C6-derivative}
\end{equation}
and dividing by $t_L(0)=2\kappa I$ gives \cref{eq:C6-hamiltonian}.
\end{proof}

As on the rational cells, the bulk sum carries a coefficient two, for the same
reason: the auxiliary line traverses every bond twice. Halving the logarithmic
derivative gives the conventional normalization, with the boundary term halved
along with it.

No analogous construction is claimed for Classes~1--5. The argument above uses
three specific facts about Class~6 --- the invertibility of the partial transpose
at nonzero rate, the polynomiality that permits continuation in the rate, and the
vanishing of both partial traces of the density --- and each has to be re-examined
before the same route is tried elsewhere.

\section{Independent checks}
\label{sec:checks}

The classification rests on the exact eliminations of
\cref{lem:stratification,thm:raw} and on the continuation argument of
\cref{subsec:continuation}. Alongside these, several independent computations
were run whose purpose is to catch an error rather than to establish a theorem,
and it is worth being explicit about which is which.

Every family printed in \cref{sec:classification} was substituted back into the
two-variable identity \cref{eq:re} in exact arithmetic, with the bulk scalars
restored and with independent symbols for the values of every arbitrary function
at $u$ and at $v$. This establishes sufficiency, which the differentiated system
does not; it would not by itself establish exhaustiveness. Across the six classes
the substitution covers $26$ cells and $11$ boundary templates for Classes~1
and~2, $10$ cells and $19$ families for Class~3 with all $23$ cell--family
incidences, $9$ cells with $13$ families and $14$ overlaps for Class~4, and the
five and four cells of Classes~5 and~6 with their orbit representatives; every
residual is identically zero.

Second, a power-series regression was run independently of the elimination. In
the gauge $k_1(u)=1$, all three remaining matrix coefficients at each order
$n=2,\dots,12$ were introduced as unrestricted variables, the reflection equation
was expanded through total degree $n+1$, the classified extension was checked
against the general solution, and the completed series was then substituted
directly through total degree $13$. The totals are those of \cref{tab:jets}: $394$
representative bulk and slope points, $4334$ extension steps, every step exact and
every completed residual zero. A rank of three at an order means the next
coefficient is unique in the fixed gauge, a rank of two means one coefficient
remains free, and a rank of zero marks the points at which a new commuting
direction can begin at higher order. These are diagnostics. They cannot exclude a
branch that first separates above order twelve, and they say nothing about global
continuation; the exact differentiated-system, functional-equation and
meromorphic-identity arguments carry those steps.

\begin{table}[t]
\centering
\caption{Independent power-series regressions through order twelve. Ranks are
those of the extension system at each order, constant over the eleven orders
tested.}
\label{tab:jets}
\begin{tabular}{c r r r r r}
\toprule
Class & Representatives & Extension steps & Rank 0 & Rank 2 & Rank 3 \\
\midrule
1 & 124 & 1364 & 8 & 71 & 45 \\
2 & 124 & 1364 & 10 & 58 & 56 \\
3 & 52 & 572 & 2 & 18 & 32 \\
4 & 67 & 737 & 6 & 26 & 35 \\
5 & 17 & 187 & 2 & 6 & 9 \\
6 & 10 & 110 & 1 & 2 & 7 \\
\midrule
Total & 394 & 4334 & 29 & 181 & 184 \\
\bottomrule
\end{tabular}
\end{table}

Third, the double-row transfer matrices of
\cref{eq:transfer,eq:C6-transfer} were built explicitly and their commutators
computed at $L=2,3,4$ and $5$, both on the rational cells of
\cref{eq:rational-cell} and, for Class~6, at a generic point $(a,Q)=(2,3)$ and at
the singular point $(0,3)$. All vanish. These finite-length checks are
regressions on the all-length arguments of
\cref{thm:rational-cell,thm:class6}, not substitutes for them; conversely the two
nonzero entries of \cref{eq:C6-no-go,eq:C5-no-go} are counterexamples, and a
counterexample at one length is conclusive in a way a confirmation at one length
is not.

Finally, the objects themselves were reconstructed independently of the
classification: the six bulk matrices and their removable limits, the Yang--Baxter
equation, regularity, braided unitarity, the crossing identity \cref{eq:crossing},
the dual identity \cref{eq:C6-dual}, the density \cref{eq:C6-density} and its two
partial traces were each verified symbolically from the definitions.

\section{Discussion and limitations}
\label{sec:discussion}

The classification splits the six families along a line that was not visible from
the bulk. Classes~1--4 share the pattern \cref{eq:pattern}, and that pattern
alone destroys crossing: their partial transposes are singular, so no crossing
datum exists at any parameter value, and the standard route to a left boundary is
closed for them entirely. Classes~5 and~6 are rational rather than exponential,
possess crossing wherever their rate is nonzero, and there the obstruction is of a
different kind --- not the absence of the datum but the failure of the left
boundary it prescribes to work against every independently chosen right boundary.
That failure is exact and quantitative, \cref{eq:C6-no-go,eq:C5-no-go}, and it
disappears precisely at the isotropic points \cref{eq:rational-cell}, where the
bulk matrix is the exchange operator plus a multiple of the identity.

What \cref{thm:class6} then shows is that the failure was an artefact of asking
for too much. An open chain needs one left boundary per right boundary, not a map
that works for every pair, and for Class~6 the constant matrix supplies it at
every bulk point. The three ingredients are the invertibility of the partial
transpose at nonzero rate, which brings the dual reflection equation into play;
polynomiality in the rate, which carries the conclusion onto the stratum where
that invertibility fails; and the vanishing of both partial traces of the
density, which is what removes the auxiliary-end term from
\cref{eq:C6-hamiltonian}. Where those three hold elsewhere, the same route should
work; the natural next case is Class~5, whose bulk matrix is the nearest
structural neighbour, and it is not settled here.

Seven limitations bound what has been established.

\begin{enumerate}
\item The load-bearing eliminations are computer assisted. Their specification is
in the text --- the objects, the test, and the coverage argument of
\cref{lem:stratification} --- so the proof can be followed without running
anything, but the ideals themselves were computed rather than derived by hand.
\item The computations have been checked internally and re-derived from the
definitions by independent programs, but no external group has reproduced them.
\item The literature audit behind the novelty statement is broad and dated, and
cannot exclude an unindexed or unpublished treatment of the same six families.
\item The power-series regressions cannot exclude a branch that first separates
above order twelve, and they do not address global continuation. Those steps rest
on the arguments of \cref{sec:exhaustiveness,subsec:continuation}.
\item On the generic Class-5 and Class-6 cells, crossing data exist but the
prescribed left map is not usable. This does not classify alternative dual
constructions outside the crossing convention \cref{eq:crossing}, and
\cref{thm:class6} is one such alternative.
\item On the rational cells, some admissible literal gauges make $t_L(0)$ zero or
singular, as \cref{eq:literal-transfer} shows; the logarithmic derivative is
asserted only in the finite nonzero case.
\item The coefficient two in \cref{eq:H-exact,eq:C6-hamiltonian} belongs to the
unhalved logarithmic derivative adopted here. It is exact, and
\cref{prop:normalization} shows that no endpoint operator removes it; the
customary half-normalization does.
\end{enumerate}

Two questions are left open and both are concrete. Does the fixed left boundary
of \cref{thm:class6} have an analogue for Classes~1--5, and if not, what replaces
it at a wall for the four classes that have no crossing datum at all? And what do
the resulting open chains look like spectrally: Classes~5 and~6 have
non-diagonalizable densities, \cref{eq:C6-density} being nilpotent in its $Q$
part, so the boundary terms of \cref{eq:C6-hamiltonian} act on a Jordan structure
rather than on a spectrum, and the effect of a boundary on that structure is not
known. There is precedent to draw on. The generalized eigenvectors of a
non-diagonalizable chain can be recovered as the limit of the eigenvectors of a
diagonalizable perturbation of it, as has been carried out for the eclectic
chain \cite{NietoGarcia_2022jb}; and the quantum-group-invariant open XXZ chain
at a root of unity is a case where an open end and Jordan cells already coexist,
and where the algebraic Bethe ansatz was extended to cover the generalized
eigenvectors \cite{Gainutdinov_2016ab}. Whether either construction survives a
boundary of the kind classified here is open.

\section{Conclusion}
\label{sec:conclusion}

For each of the six regular difference-form $4\times4$ bulk $R$-matrices outside
the eight-vertex list, and at every point of its full complex parameter variety,
we have determined every regular single-valued meromorphic right reflection
matrix, together with its orbit under the scalar gauge and the constant
similarities and spectral rescalings that preserve the same literal bulk matrix.
The derivation assumes no form for the boundary: one differentiation of the
reflection equation at the origin makes the problem linear, and an exact
stratification by the rank of the resulting system splits each parameter variety
into finitely many cells carrying either a unique kernel or a centralizer family
of arbitrary functions. Every removable limit, singular stratum,
arbitrary-function fiber and overlap is included, and every family is a
single-valued meromorphic function on the plane.

The consequences for open chains are sharp and asymmetric. Four of the six
classes admit no crossing datum at any parameter value, for the structural reason
that their first-leg partial transposes drop rank; the other two admit a unique
datum exactly at nonzero rate, and the left boundary it prescribes satisfies the
every-pair requirement only at their isotropic points. Fixing the left boundary to
the identity instead gives, for Class~6 and at every bulk point, a commuting
double-row transfer matrix for every classified right boundary, whose logarithmic
derivative is twice the bulk chain plus one one-site endpoint term.

What this opens is no longer a classification problem. The exact boundary
families make the spectral, Jordan-block and continuum-limit analyses of these
open chains concrete, and the six bulk models --- several of them non-Hermitian
and non-diagonalizable --- can now be studied with a wall. The first thing to
settle is whether the constant left boundary of \cref{thm:class6} has an analogue
in the remaining five classes; a single Class-5 counterexample, at any bulk point
and any chain length, would decide it in the negative.

\section*{Acknowledgments}

Z.\ Y.\ acknowledges support from the Strategic Priority Research Program of the
Chinese Academy of Sciences under Grant No.~XDB1680102. Z.\ Y.\ also
acknowledges support by National Natural Science Foundation of China (Grant
No.~12247101). We also acknowledge Gewu Intelligence Lab for providing the
collaborative research environment in which this project was developed.

OpenAI GPT-5.6 Sol was used substantively in this work: in solving the
problem, including the formulation of arguments and the derivation of
intermediate steps, and in drafting and revising the manuscript. The authors
specified the problem, guided the approach to be taken, checked every argument
and result it produced, and take full responsibility for the content of this
manuscript.

\section*{Code availability}

Four steps of the argument are computer assisted: the constructible
stratification of \cref{lem:stratification}, on which \cref{thm:raw} rests; the
tensor-square generator computation behind \cref{thm:quotient}; the coefficient
elimination establishing \cref{prop:crossing-56}; and the substitution verifying
that the constant left boundary satisfies the dual reflection equation
\cref{eq:C6-dual}, which is the hypothesis discharged in the proof of
\cref{thm:class6}. Each is specified in the text --- the objects, the test applied
to them, and why the cases are exhaustive --- so that the proof can be followed
without running anything. The exact symbolic programs that carry out these
computations, and the programs performing the independent checks of
\cref{sec:checks}, are available from the author on request.

\bibliographystyle{jphysa}
\bibliography{references,unverified}

\end{document}
